\documentclass[../main.tex]{subfiles}
\begin{document}

\chapter{Parallel Programming}
\lessoninfo{
  video={M2, lessons 1--4 (Parallel Programming Patterns; OpenMP vs.\ MPI, parts 1--2; Programming with MPI)},
  sources={\cite{barney2026parallel} (required); \cite{openmp2021spec}; \cite{mpi2023standard}},
  textbook={PPP \cite{mattson2004patterns} ch.~3--5, App.~A--B (optional); PMPP \cite{hwu2022pmpp} ch.~20},
  keywords={parallel programming pattern, master/worker, loop parallelism, fork/join, pipeline, shared memory, distributed memory, OpenMP, thread, reduction, data race, critical section, scheduling, barrier, MPI, stencil, halo},
}

Lesson~1 introduced the GPU as a throughput machine and the two ways of
splitting a program into parallel work, task decomposition and data
decomposition.  This lesson shows how that work is expressed in code on a
conventional multiprocessor, before the GPU-specific model of Lesson~3.
It first collects the recurring \emph{patterns} of parallel programs, then
contrasts the two memory models a parallel program can be written against:
shared memory, where threads communicate through one address space and are
programmed with OpenMP, and distributed memory, where processes exchange
messages with MPI.  Both matter for CUDA, which combines them: some memory
is shared by every thread of a block, while other data must be moved
explicitly.  The chapter closes with a stencil computation whose boundary
exchange between processes returns, on the GPU, as the halo loading of
Lesson~3.

\section{Parallel programming patterns}
\label{sec:l02-patterns}

A parallel application is built from a serial one in three moves: the work
is divided into pieces, the data structures are adjusted so that the pieces
can be worked on independently, and the pieces are distributed to the
processors (\cref{fig:l02-parallel-app}).\source{M2, parallel programming
patterns; PPP ch.~3--5.}  A \term{parallel programming pattern}[並列プログラミングパターン]
is a reusable way of organising this division and distribution; the handful
of patterns below cover most parallel programs, including the GPU kernels
of the later lessons.

\begin{figure}[tb]
  \centering
  % Creating a parallel application in three steps: the original program's
% work is divided and its data structures adjusted, then the pieces are
% distributed to processors.  Redraws lowyx m2l1_parallelapp.
% Usage: % Creating a parallel application in three steps: the original program's
% work is divided and its data structures adjusted, then the pieces are
% distributed to processors.  Redraws lowyx m2l1_parallelapp.
% Usage: % Creating a parallel application in three steps: the original program's
% work is divided and its data structures adjusted, then the pieces are
% distributed to processors.  Redraws lowyx m2l1_parallelapp.
% Usage: \input{figures/l02-parallel-app}   (body width, 120 mm)
\begin{tikzpicture}[x=1mm, y=1mm, font=\scriptsize\sffamily,
                    >={Stealth[length=1.4mm]}]
  % panel 1: original program
  \node[anchor=south, font=\scriptsize\sffamily\bfseries, align=center]
    at (12,21) {\iflangja{元のプログラム}{Original program}};
  \fill[yellow!45, draw=ln-muted] (3,0) rectangle (21,18);
  \fill[orange!55, draw=ln-muted] (3,-3) rectangle (21,-1);
  \fill[green!45, draw=ln-muted] (3,-6) rectangle (21,-4);
  \draw[->, thick] (24,8) -- (31,8);
  % panel 2: divided work
  \node[anchor=south, font=\scriptsize\sffamily\bfseries, align=center]
    at (47.5,21) {\iflangja{仕事を分割し\\データ構造を調整}{Divide the work,\\adjust the data structures}};
  \foreach \i in {0,...,5} {
    \fill[yellow!45, draw=ln-muted] (35+\i*4.3,0) rectangle (38.5+\i*4.3,18);
  }
  \fill[orange!55, draw=ln-muted] (35,-3) rectangle (47,-1);
  \fill[orange!55, draw=ln-muted] (48,-3) rectangle (60,-1);
  \fill[green!45, draw=ln-muted] (35,-6) rectangle (60,-4);
  \draw[->, thick] (63,8) -- (70,8);
  % panel 3: processors
  \node[anchor=south, font=\scriptsize\sffamily\bfseries, align=center]
    at (93,21) {\iflangja{仕事を配る}{Distribute the work}};
  \foreach \i in {0,...,5} {
    \draw[ln-accent, thick] (74+\i*6.6,2) rectangle (79.4+\i*6.6,14);
    \node at (76.7+\i*6.6,8) {P\i};
  }
  \node[text=ln-muted] at (93,-2) {\iflangja{プロセッサ}{processors}};
\end{tikzpicture}
   (body width, 120 mm)
\begin{tikzpicture}[x=1mm, y=1mm, font=\scriptsize\sffamily,
                    >={Stealth[length=1.4mm]}]
  % panel 1: original program
  \node[anchor=south, font=\scriptsize\sffamily\bfseries, align=center]
    at (12,21) {\iflangja{元のプログラム}{Original program}};
  \fill[yellow!45, draw=ln-muted] (3,0) rectangle (21,18);
  \fill[orange!55, draw=ln-muted] (3,-3) rectangle (21,-1);
  \fill[green!45, draw=ln-muted] (3,-6) rectangle (21,-4);
  \draw[->, thick] (24,8) -- (31,8);
  % panel 2: divided work
  \node[anchor=south, font=\scriptsize\sffamily\bfseries, align=center]
    at (47.5,21) {\iflangja{仕事を分割し\\データ構造を調整}{Divide the work,\\adjust the data structures}};
  \foreach \i in {0,...,5} {
    \fill[yellow!45, draw=ln-muted] (35+\i*4.3,0) rectangle (38.5+\i*4.3,18);
  }
  \fill[orange!55, draw=ln-muted] (35,-3) rectangle (47,-1);
  \fill[orange!55, draw=ln-muted] (48,-3) rectangle (60,-1);
  \fill[green!45, draw=ln-muted] (35,-6) rectangle (60,-4);
  \draw[->, thick] (63,8) -- (70,8);
  % panel 3: processors
  \node[anchor=south, font=\scriptsize\sffamily\bfseries, align=center]
    at (93,21) {\iflangja{仕事を配る}{Distribute the work}};
  \foreach \i in {0,...,5} {
    \draw[ln-accent, thick] (74+\i*6.6,2) rectangle (79.4+\i*6.6,14);
    \node at (76.7+\i*6.6,8) {P\i};
  }
  \node[text=ln-muted] at (93,-2) {\iflangja{プロセッサ}{processors}};
\end{tikzpicture}
   (body width, 120 mm)
\begin{tikzpicture}[x=1mm, y=1mm, font=\scriptsize\sffamily,
                    >={Stealth[length=1.4mm]}]
  % panel 1: original program
  \node[anchor=south, font=\scriptsize\sffamily\bfseries, align=center]
    at (12,21) {\iflangja{元のプログラム}{Original program}};
  \fill[yellow!45, draw=ln-muted] (3,0) rectangle (21,18);
  \fill[orange!55, draw=ln-muted] (3,-3) rectangle (21,-1);
  \fill[green!45, draw=ln-muted] (3,-6) rectangle (21,-4);
  \draw[->, thick] (24,8) -- (31,8);
  % panel 2: divided work
  \node[anchor=south, font=\scriptsize\sffamily\bfseries, align=center]
    at (47.5,21) {\iflangja{仕事を分割し\\データ構造を調整}{Divide the work,\\adjust the data structures}};
  \foreach \i in {0,...,5} {
    \fill[yellow!45, draw=ln-muted] (35+\i*4.3,0) rectangle (38.5+\i*4.3,18);
  }
  \fill[orange!55, draw=ln-muted] (35,-3) rectangle (47,-1);
  \fill[orange!55, draw=ln-muted] (48,-3) rectangle (60,-1);
  \fill[green!45, draw=ln-muted] (35,-6) rectangle (60,-4);
  \draw[->, thick] (63,8) -- (70,8);
  % panel 3: processors
  \node[anchor=south, font=\scriptsize\sffamily\bfseries, align=center]
    at (93,21) {\iflangja{仕事を配る}{Distribute the work}};
  \foreach \i in {0,...,5} {
    \draw[ln-accent, thick] (74+\i*6.6,2) rectangle (79.4+\i*6.6,14);
    \node at (76.7+\i*6.6,8) {P\i};
  }
  \node[text=ln-muted] at (93,-2) {\iflangja{プロセッサ}{processors}};
\end{tikzpicture}

  \caption{Creating a parallel application: divide the work, adjust the
    data structures, distribute the pieces.  The orange and green bars are
    parts of the program that split differently from the main body.}
  \label{fig:l02-parallel-app}
\end{figure}

\subsection{Four steps}

Designing a parallel program proceeds in four steps.\sidenote{The steps
follow the four \emph{design spaces} of \cite{mattson2004patterns}: Finding
Concurrency, Algorithm Structure, Supporting Structures and Implementation
Mechanisms.  The book places the patterns below in different spaces
(Pipeline in Algorithm Structure; SPMD, Master/Worker, Loop Parallelism and
Fork/Join in Supporting Structures); the course lists them side by side.}
\begin{enumerate}
  \item \textbf{Discover concurrency.}  The \term{concurrency}[並行性] of a
    computation is the extent to which its parts may proceed independently
    of one another; the first step finds where such independence exists
    (which tasks, which data).
  \item \textbf{Structure the algorithm} so that the concurrency can be
    exploited: decide what the units of parallel work are and how they
    communicate.
  \item \textbf{Implement} it in a language and toolset suited to the
    target: OpenMP, MPI, CUDA.
  \item \textbf{Execute and optimise} on the parallel system, measuring
    and tuning until the performance goal is met.
\end{enumerate}

\subsection{Five patterns}

\begin{description}
  \item[Master/worker]  In the \term{master/worker pattern}[マスタ／ワーカパターン]
    one master thread or process sets up the problem, places the pieces of
    work in a \term{task queue}[タスクキュー] (a \enquote{bag of tasks}),
    launches a pool of workers, waits, collects their results and
    terminates the computation (\cref{fig:l02-master-worker}).  Each worker
    repeatedly dequeues a task and runs it; since workers take the next
    task as soon as they finish the previous one, the load balances itself.
    The pattern suits work that splits into independent pieces of unequal
    or unknown size.\margin{The OpenMP \texttt{dynamic} schedule of
    \cref{sec:l02-scheduling} is a task queue whose tasks are chunks of
    loop iterations.}
  \item[Loop parallelism]  In \term{loop parallelism}[ループ並列性] the
    iterations of a loop run concurrently when they are independent, as in
\begin{lstlisting}[language=C++, numbers=none]
for (int i = 0; i < 16; i++) c[i] = a[i] + b[i];
\end{lstlisting}
    Unlike master/worker tasks, the iterations are usually identical
    operations on different data, which makes this the pattern of
    \term*{data decomposition} (Lesson~1).\margin{The same vector addition is
    the first CUDA kernel of Lesson~3: one thread per iteration.}
  \item[SPMD]  Every processing element runs the same program on its own
    portion of the data; \term*{SPMD} (Lesson~1) is the usual choice for GPU
    programming, where thousands of threads run one kernel and pick their
    data by their thread index.
  \item[Fork/join]  In the \term{fork/join pattern}[フォーク／ジョインパターン]
    a parent task forks child tasks, continues or waits, and joins them
    (waits for their completion) before going on
    (\cref{fig:l02-fork-join}).  It mixes serial and parallel phases from
    a single entry point; an OpenMP parallel region is a fork/join.
  \item[Pipeline]  The \term{pipeline pattern}[パイプラインパターン]
    assigns each processor one stage of a multi-stage computation and
    streams the data items through the stages, exactly as a CPU pipeline
    streams instructions (\cref{fig:l02-pipeline}).  It suits data
    streams: signal processing, the graphics pipeline of Lesson~1, or a
    compression workflow (decompress, process, compress).
\end{description}

\begin{figure}[tb]
  \centering
  % The master/worker pattern: a master sets up the problem, fills a task
% queue (bag of tasks), launches workers, waits for them and collects the
% results.  Workers dequeue tasks until the queue is empty.
% Redraws lowyx m2l1_master_worker_pattern.
% Usage: % The master/worker pattern: a master sets up the problem, fills a task
% queue (bag of tasks), launches workers, waits for them and collects the
% results.  Workers dequeue tasks until the queue is empty.
% Redraws lowyx m2l1_master_worker_pattern.
% Usage: % The master/worker pattern: a master sets up the problem, fills a task
% queue (bag of tasks), launches workers, waits for them and collects the
% results.  Workers dequeue tasks until the queue is empty.
% Redraws lowyx m2l1_master_worker_pattern.
% Usage: \input{figures/l02-master-worker}   (body width, 120 mm)
\begin{tikzpicture}[x=1mm, y=1mm, font=\scriptsize\sffamily,
                    >={Stealth[length=1.4mm]}]
  \node[draw=ln-accent, thick, align=left, anchor=west, inner sep=3pt,
        text width=31mm] (m) at (0,8)
    {\textbf{\iflangja{マスタ}{Master}}\\[1pt]
     1\ \iflangja{計算を開始}{initiate the computation}\\
     2\ \iflangja{問題を設定}{set up the problem}\\
     3\ \iflangja{タスクの袋を作る}{create the bag of tasks}\\
     4\ \iflangja{ワーカを起動}{launch the workers}\\
     5\ \iflangja{待つ}{wait}\\
     6\ \iflangja{結果を集める}{collect the results}\\
     7\ \iflangja{計算を終了}{terminate}};
  % task queue
  \node[draw=ln-muted, dashed, inner sep=0pt, minimum width=30mm,
        minimum height=11mm] (q) at (58,12) {};
  \node[above=1pt of q, text=ln-muted] {\iflangja{タスクキュー（タスクの袋）}{task queue (bag of tasks)}};
  \foreach \i/\c in {0/yellow!60, 1/orange!60, 2/green!50, 3/blue!30, 4/red!35, 5/purple!30} {
    \fill[\c, draw=ln-muted] (45+\i*4.6,8.5) rectangle (48.6+\i*4.6,15.5);
  }
  \node[font=\scriptsize\sffamily] at (46.8,12) {t};
  % workers
  \foreach \i in {0,1,2} {
    \node[draw=ln-accent, thick, minimum width=14mm, minimum height=6mm]
      (w\i) at (98,20-\i*8) {\iflangja{ワーカ}{Worker} \i};
    \draw[->] (q.east) -- (w\i.west);
  }
  \node[text=ln-muted, anchor=west] at (76,24) {\iflangja{デキュー}{dequeue}};
  \draw[->] (m.east |- q.west) -- (q.west) node[midway, above] {\iflangja{タスク}{tasks}};
  % results back to the master
  \draw[->, ln-accent] (w2.south) |- node[pos=0.25, right, text=ln-accent] {\iflangja{結果}{results}}
    ([yshift=-3mm]m.south) -- (m.south);
\end{tikzpicture}
   (body width, 120 mm)
\begin{tikzpicture}[x=1mm, y=1mm, font=\scriptsize\sffamily,
                    >={Stealth[length=1.4mm]}]
  \node[draw=ln-accent, thick, align=left, anchor=west, inner sep=3pt,
        text width=31mm] (m) at (0,8)
    {\textbf{\iflangja{マスタ}{Master}}\\[1pt]
     1\ \iflangja{計算を開始}{initiate the computation}\\
     2\ \iflangja{問題を設定}{set up the problem}\\
     3\ \iflangja{タスクの袋を作る}{create the bag of tasks}\\
     4\ \iflangja{ワーカを起動}{launch the workers}\\
     5\ \iflangja{待つ}{wait}\\
     6\ \iflangja{結果を集める}{collect the results}\\
     7\ \iflangja{計算を終了}{terminate}};
  % task queue
  \node[draw=ln-muted, dashed, inner sep=0pt, minimum width=30mm,
        minimum height=11mm] (q) at (58,12) {};
  \node[above=1pt of q, text=ln-muted] {\iflangja{タスクキュー（タスクの袋）}{task queue (bag of tasks)}};
  \foreach \i/\c in {0/yellow!60, 1/orange!60, 2/green!50, 3/blue!30, 4/red!35, 5/purple!30} {
    \fill[\c, draw=ln-muted] (45+\i*4.6,8.5) rectangle (48.6+\i*4.6,15.5);
  }
  \node[font=\scriptsize\sffamily] at (46.8,12) {t};
  % workers
  \foreach \i in {0,1,2} {
    \node[draw=ln-accent, thick, minimum width=14mm, minimum height=6mm]
      (w\i) at (98,20-\i*8) {\iflangja{ワーカ}{Worker} \i};
    \draw[->] (q.east) -- (w\i.west);
  }
  \node[text=ln-muted, anchor=west] at (76,24) {\iflangja{デキュー}{dequeue}};
  \draw[->] (m.east |- q.west) -- (q.west) node[midway, above] {\iflangja{タスク}{tasks}};
  % results back to the master
  \draw[->, ln-accent] (w2.south) |- node[pos=0.25, right, text=ln-accent] {\iflangja{結果}{results}}
    ([yshift=-3mm]m.south) -- (m.south);
\end{tikzpicture}
   (body width, 120 mm)
\begin{tikzpicture}[x=1mm, y=1mm, font=\scriptsize\sffamily,
                    >={Stealth[length=1.4mm]}]
  \node[draw=ln-accent, thick, align=left, anchor=west, inner sep=3pt,
        text width=31mm] (m) at (0,8)
    {\textbf{\iflangja{マスタ}{Master}}\\[1pt]
     1\ \iflangja{計算を開始}{initiate the computation}\\
     2\ \iflangja{問題を設定}{set up the problem}\\
     3\ \iflangja{タスクの袋を作る}{create the bag of tasks}\\
     4\ \iflangja{ワーカを起動}{launch the workers}\\
     5\ \iflangja{待つ}{wait}\\
     6\ \iflangja{結果を集める}{collect the results}\\
     7\ \iflangja{計算を終了}{terminate}};
  % task queue
  \node[draw=ln-muted, dashed, inner sep=0pt, minimum width=30mm,
        minimum height=11mm] (q) at (58,12) {};
  \node[above=1pt of q, text=ln-muted] {\iflangja{タスクキュー（タスクの袋）}{task queue (bag of tasks)}};
  \foreach \i/\c in {0/yellow!60, 1/orange!60, 2/green!50, 3/blue!30, 4/red!35, 5/purple!30} {
    \fill[\c, draw=ln-muted] (45+\i*4.6,8.5) rectangle (48.6+\i*4.6,15.5);
  }
  \node[font=\scriptsize\sffamily] at (46.8,12) {t};
  % workers
  \foreach \i in {0,1,2} {
    \node[draw=ln-accent, thick, minimum width=14mm, minimum height=6mm]
      (w\i) at (98,20-\i*8) {\iflangja{ワーカ}{Worker} \i};
    \draw[->] (q.east) -- (w\i.west);
  }
  \node[text=ln-muted, anchor=west] at (76,24) {\iflangja{デキュー}{dequeue}};
  \draw[->] (m.east |- q.west) -- (q.west) node[midway, above] {\iflangja{タスク}{tasks}};
  % results back to the master
  \draw[->, ln-accent] (w2.south) |- node[pos=0.25, right, text=ln-accent] {\iflangja{結果}{results}}
    ([yshift=-3mm]m.south) -- (m.south);
\end{tikzpicture}

  \caption{The master/worker pattern.  The master fills the task queue
    and waits; workers dequeue tasks until the queue is empty and return
    their results.}
  \label{fig:l02-master-worker}
\end{figure}

\begin{marginfigure}
  \resizebox{\linewidth}{!}{% The fork/join pattern: a single flow of control forks into parallel
% tasks inside a parallel region and joins back to one flow.
% Redraws lowyx m2l1_forkjoin_pattern.
% Usage: % The fork/join pattern: a single flow of control forks into parallel
% tasks inside a parallel region and joins back to one flow.
% Redraws lowyx m2l1_forkjoin_pattern.
% Usage: % The fork/join pattern: a single flow of control forks into parallel
% tasks inside a parallel region and joins back to one flow.
% Redraws lowyx m2l1_forkjoin_pattern.
% Usage: \input{figures/l02-fork-join}   (margin width, 50 mm)
\begin{tikzpicture}[x=1mm, y=1mm, font=\scriptsize\sffamily,
                    >={Stealth[length=1.4mm]}]
  \fill[yellow!35, draw=ln-muted, dashed] (14,-9) rectangle (36,9);
  \node[anchor=south, text=ln-muted] at (25,9.5) {\iflangja{並列領域}{parallel region}};
  \draw[thick, ln-accent] (0,0) -- (14,0);
  \foreach \y in {-6,-2,2,6} {
    \draw[thick, ln-accent] (14,0) -- (18,\y) -- (32,\y) -- (36,0);
  }
  \draw[->, thick, ln-accent] (36,0) -- (50,0);
  \node[below, text=ln-muted] at (14,-9.5) {\iflangja{フォーク}{fork}};
  \node[below, text=ln-muted] at (36,-9.5) {\iflangja{ジョイン}{join}};
  \node[above, text=ln-muted] at (7,0.5) {\iflangja{逐次}{serial}};
  \node[above, text=ln-muted] at (43,0.5) {\iflangja{逐次}{serial}};
\end{tikzpicture}
   (margin width, 50 mm)
\begin{tikzpicture}[x=1mm, y=1mm, font=\scriptsize\sffamily,
                    >={Stealth[length=1.4mm]}]
  \fill[yellow!35, draw=ln-muted, dashed] (14,-9) rectangle (36,9);
  \node[anchor=south, text=ln-muted] at (25,9.5) {\iflangja{並列領域}{parallel region}};
  \draw[thick, ln-accent] (0,0) -- (14,0);
  \foreach \y in {-6,-2,2,6} {
    \draw[thick, ln-accent] (14,0) -- (18,\y) -- (32,\y) -- (36,0);
  }
  \draw[->, thick, ln-accent] (36,0) -- (50,0);
  \node[below, text=ln-muted] at (14,-9.5) {\iflangja{フォーク}{fork}};
  \node[below, text=ln-muted] at (36,-9.5) {\iflangja{ジョイン}{join}};
  \node[above, text=ln-muted] at (7,0.5) {\iflangja{逐次}{serial}};
  \node[above, text=ln-muted] at (43,0.5) {\iflangja{逐次}{serial}};
\end{tikzpicture}
   (margin width, 50 mm)
\begin{tikzpicture}[x=1mm, y=1mm, font=\scriptsize\sffamily,
                    >={Stealth[length=1.4mm]}]
  \fill[yellow!35, draw=ln-muted, dashed] (14,-9) rectangle (36,9);
  \node[anchor=south, text=ln-muted] at (25,9.5) {\iflangja{並列領域}{parallel region}};
  \draw[thick, ln-accent] (0,0) -- (14,0);
  \foreach \y in {-6,-2,2,6} {
    \draw[thick, ln-accent] (14,0) -- (18,\y) -- (32,\y) -- (36,0);
  }
  \draw[->, thick, ln-accent] (36,0) -- (50,0);
  \node[below, text=ln-muted] at (14,-9.5) {\iflangja{フォーク}{fork}};
  \node[below, text=ln-muted] at (36,-9.5) {\iflangja{ジョイン}{join}};
  \node[above, text=ln-muted] at (7,0.5) {\iflangja{逐次}{serial}};
  \node[above, text=ln-muted] at (43,0.5) {\iflangja{逐次}{serial}};
\end{tikzpicture}
}
  \caption{The fork/join pattern: one flow of control forks into four
    parallel tasks and joins again.}
  \label{fig:l02-fork-join}
\end{marginfigure}

\begin{figure}[tb]
  \centering
  % The pipeline pattern: four work items pass through five stages; at any
% time step up to four items are in different stages, like instructions in
% a CPU pipeline.  Redraws lowyx m2l1_pipeline_pattern.
% Usage: % The pipeline pattern: four work items pass through five stages; at any
% time step up to four items are in different stages, like instructions in
% a CPU pipeline.  Redraws lowyx m2l1_pipeline_pattern.
% Usage: % The pipeline pattern: four work items pass through five stages; at any
% time step up to four items are in different stages, like instructions in
% a CPU pipeline.  Redraws lowyx m2l1_pipeline_pattern.
% Usage: \input{figures/l02-pipeline}   (body width, 120 mm)
\begin{tikzpicture}[x=1mm, y=1mm, font=\scriptsize\sffamily,
                    >={Stealth[length=1.4mm]}]
  \def\cw{8}  % cell width
  \def\ch{6}  % cell height
  % time axis
  \draw[->] (16,1) -- (16+9*\cw,1) node[below left] {\iflangja{時間ステップ}{time step}};
  \foreach \t in {1,...,8} {
    \node[above, text=ln-muted] at (16+\t*\cw-\cw/2,1) {\t};
  }
  % rows: work items; each row shifted one step right
  \foreach \r in {0,1,2,3} {
    \node[anchor=east] at (15,-3-\r*\ch) {\iflangja{仕事}{Work} \the\numexpr\r+1\relax};
    \foreach \s/\c in {0/red!40, 1/green!45, 2/purple!35, 3/yellow!60, 4/gray!40} {
      \fill[\c, draw=ln-muted] (16+\r*\cw+\s*\cw,-\ch-\r*\ch)
        rectangle (16+\r*\cw+\s*\cw+\cw,-\r*\ch);
      \node at (16+\r*\cw+\s*\cw+\cw/2,-\ch/2-\r*\ch) {S\the\numexpr\s+1\relax};
    }
  }
  % legend
  \node[anchor=west, text=ln-muted, align=left] at (16,-31)
    {S1\,--\,S5: \iflangja{パイプラインの段}{pipeline stages}
     (\iflangja{例：伸長，処理，圧縮}{e.g.\ decompress, process, compress})};
\end{tikzpicture}
   (body width, 120 mm)
\begin{tikzpicture}[x=1mm, y=1mm, font=\scriptsize\sffamily,
                    >={Stealth[length=1.4mm]}]
  \def\cw{8}  % cell width
  \def\ch{6}  % cell height
  % time axis
  \draw[->] (16,1) -- (16+9*\cw,1) node[below left] {\iflangja{時間ステップ}{time step}};
  \foreach \t in {1,...,8} {
    \node[above, text=ln-muted] at (16+\t*\cw-\cw/2,1) {\t};
  }
  % rows: work items; each row shifted one step right
  \foreach \r in {0,1,2,3} {
    \node[anchor=east] at (15,-3-\r*\ch) {\iflangja{仕事}{Work} \the\numexpr\r+1\relax};
    \foreach \s/\c in {0/red!40, 1/green!45, 2/purple!35, 3/yellow!60, 4/gray!40} {
      \fill[\c, draw=ln-muted] (16+\r*\cw+\s*\cw,-\ch-\r*\ch)
        rectangle (16+\r*\cw+\s*\cw+\cw,-\r*\ch);
      \node at (16+\r*\cw+\s*\cw+\cw/2,-\ch/2-\r*\ch) {S\the\numexpr\s+1\relax};
    }
  }
  % legend
  \node[anchor=west, text=ln-muted, align=left] at (16,-31)
    {S1\,--\,S5: \iflangja{パイプラインの段}{pipeline stages}
     (\iflangja{例：伸長，処理，圧縮}{e.g.\ decompress, process, compress})};
\end{tikzpicture}
   (body width, 120 mm)
\begin{tikzpicture}[x=1mm, y=1mm, font=\scriptsize\sffamily,
                    >={Stealth[length=1.4mm]}]
  \def\cw{8}  % cell width
  \def\ch{6}  % cell height
  % time axis
  \draw[->] (16,1) -- (16+9*\cw,1) node[below left] {\iflangja{時間ステップ}{time step}};
  \foreach \t in {1,...,8} {
    \node[above, text=ln-muted] at (16+\t*\cw-\cw/2,1) {\t};
  }
  % rows: work items; each row shifted one step right
  \foreach \r in {0,1,2,3} {
    \node[anchor=east] at (15,-3-\r*\ch) {\iflangja{仕事}{Work} \the\numexpr\r+1\relax};
    \foreach \s/\c in {0/red!40, 1/green!45, 2/purple!35, 3/yellow!60, 4/gray!40} {
      \fill[\c, draw=ln-muted] (16+\r*\cw+\s*\cw,-\ch-\r*\ch)
        rectangle (16+\r*\cw+\s*\cw+\cw,-\r*\ch);
      \node at (16+\r*\cw+\s*\cw+\cw/2,-\ch/2-\r*\ch) {S\the\numexpr\s+1\relax};
    }
  }
  % legend
  \node[anchor=west, text=ln-muted, align=left] at (16,-31)
    {S1\,--\,S5: \iflangja{パイプラインの段}{pipeline stages}
     (\iflangja{例：伸長，処理，圧縮}{e.g.\ decompress, process, compress})};
\end{tikzpicture}

  \caption{The pipeline pattern.  Four work items pass through five
    stages; at steps 4 and~5 four of the five stages are busy with
    different items (with five or more items in flight every stage would
    be).  With $s$ stages and $n$ items the pipeline takes $s+n-1$ steps
    instead of $sn$: here 8 instead of 20.}
  \label{fig:l02-pipeline}
\end{figure}

\begin{keypoint}
  Patterns are not exclusive: a GPU program is SPMD at the level of the
  kernel, loop parallelism inside it (one thread per iteration), and
  fork/join at the level of the host that launches the kernel and waits
  for it.
\end{keypoint}

\section{Shared-memory programming with OpenMP}
\label{sec:l02-openmp}

How the pieces of a parallel program communicate depends on whether they
share a memory.\source{M2, OpenMP vs.\ MPI (part 1); PPP App.~A;
\cite{barney2026parallel}, memory architectures and programming models.}

\subsection{Shared versus distributed memory}

\begin{definition}[Shared-memory programming]
  In \term{shared-memory programming}[共有メモリプログラミング] all
  processors read and write one memory: a value that processor~1 writes to
  an address is seen by processor~2 when it reads that address
  (\cref{fig:l02-memory-models}, left).  Communication is implicit in
  ordinary loads and stores.
\end{definition}

\begin{definition}[Distributed-memory programming]
  In \term{distributed-memory programming}[分散メモリプログラミング] each
  processor has its own memory, invisible to the others.  Data that lives
  in another processor's memory must be requested and sent in a
  \emph{message} (\cref{fig:l02-memory-models}, right): processor~2 asks
  processors 1 and 3 for the contents of an address in their memories,
  adds the three values and writes the result back into its own memory.
\end{definition}

\begin{figure}[tb]
  \centering
  % Shared memory (left): every processor reads and writes one memory, so a
% value written by P1 is read by P2.  Distributed memory (right): each
% processor owns a memory; P2 obtains the values of the same address in
% the other memories through messages and adds them.
% Redraws lowyx m2l2_shared_mem and m2l2_shared_mem_distributed.
% Usage: % Shared memory (left): every processor reads and writes one memory, so a
% value written by P1 is read by P2.  Distributed memory (right): each
% processor owns a memory; P2 obtains the values of the same address in
% the other memories through messages and adds them.
% Redraws lowyx m2l2_shared_mem and m2l2_shared_mem_distributed.
% Usage: % Shared memory (left): every processor reads and writes one memory, so a
% value written by P1 is read by P2.  Distributed memory (right): each
% processor owns a memory; P2 obtains the values of the same address in
% the other memories through messages and adds them.
% Redraws lowyx m2l2_shared_mem and m2l2_shared_mem_distributed.
% Usage: \input{figures/l02-memory-models}   (body width, 120 mm)
\begin{tikzpicture}[x=1mm, y=1mm, font=\scriptsize\sffamily,
                    >={Stealth[length=1.4mm]},
                    proc/.style={draw=ln-accent, thick, minimum width=13mm, minimum height=6mm},
                    mem/.style={draw=ln-muted, fill=ln-box, minimum height=7mm, align=center}]
  % ---- shared memory
  \node[anchor=west, font=\scriptsize\sffamily\bfseries] at (0,24)
    {\iflangja{共有メモリ}{Shared memory}};
  \foreach \i in {1,2,3} {
    \node[proc] (p\i) at ({8+(\i-1)*16},16) {P\i};
  }
  \node[mem, minimum width=42mm] (sm) at (24,2)
    {\iflangja{メモリ}{Memory}\\\texttt{mem[0x100] = 10}};
  \draw[->] (p1.south) -- (sm.north -| p1.south)
    node[midway, left, align=right, text=ln-muted] {\iflangja{10を\\書く}{writes 10}};
  \draw[<-] (p2.south) -- (sm.north -| p2.south)
    node[midway, right, text=ln-muted] {\iflangja{10を読む}{reads 10}};
  \draw[<->] (p3.south) -- (sm.north -| p3.south);
  % ---- distributed memory
  \node[anchor=west, font=\scriptsize\sffamily\bfseries] at (54,24)
    {\iflangja{分散メモリ}{Distributed memory}};
  \foreach \i/\v/\c in {1/3/orange!45, 2/5/blue!25, 3/8/red!35} {
    \node[proc] (q\i) at ({62+(\i-1)*22},16) {P\i};
    \node[mem, fill=\c, minimum width=17mm] (dm\i) at ({62+(\i-1)*22},2)
      {\iflangja{メモリ}{Memory} \i\\\texttt{mem[0x100] = \v}};
    \draw[<->] (q\i.south) -- (dm\i.north);
  }
  % messages to P2
  \draw[->, ln-accent, thick] (dm1.north east) to[bend left=15] (q2.west);
  \draw[->, ln-accent, thick] (dm3.north west) to[bend right=15] (q2.east);
  \node[draw=ln-accent, fill=white, inner sep=1pt, font=\scriptsize\sffamily, text=ln-accent] at (73.5,10.5) {$\boxtimes$};
  \node[draw=ln-accent, fill=white, inner sep=1pt, font=\scriptsize\sffamily, text=ln-accent] at (94.5,10.5) {$\boxtimes$};
  \node[above=1pt of q2, text=ln-accent] {$3+5+8=16$};
  \node[anchor=west, text=ln-muted] at (54,-6) {$\boxtimes$ = \iflangja{メッセージ}{message}};
\end{tikzpicture}
   (body width, 120 mm)
\begin{tikzpicture}[x=1mm, y=1mm, font=\scriptsize\sffamily,
                    >={Stealth[length=1.4mm]},
                    proc/.style={draw=ln-accent, thick, minimum width=13mm, minimum height=6mm},
                    mem/.style={draw=ln-muted, fill=ln-box, minimum height=7mm, align=center}]
  % ---- shared memory
  \node[anchor=west, font=\scriptsize\sffamily\bfseries] at (0,24)
    {\iflangja{共有メモリ}{Shared memory}};
  \foreach \i in {1,2,3} {
    \node[proc] (p\i) at ({8+(\i-1)*16},16) {P\i};
  }
  \node[mem, minimum width=42mm] (sm) at (24,2)
    {\iflangja{メモリ}{Memory}\\\texttt{mem[0x100] = 10}};
  \draw[->] (p1.south) -- (sm.north -| p1.south)
    node[midway, left, align=right, text=ln-muted] {\iflangja{10を\\書く}{writes 10}};
  \draw[<-] (p2.south) -- (sm.north -| p2.south)
    node[midway, right, text=ln-muted] {\iflangja{10を読む}{reads 10}};
  \draw[<->] (p3.south) -- (sm.north -| p3.south);
  % ---- distributed memory
  \node[anchor=west, font=\scriptsize\sffamily\bfseries] at (54,24)
    {\iflangja{分散メモリ}{Distributed memory}};
  \foreach \i/\v/\c in {1/3/orange!45, 2/5/blue!25, 3/8/red!35} {
    \node[proc] (q\i) at ({62+(\i-1)*22},16) {P\i};
    \node[mem, fill=\c, minimum width=17mm] (dm\i) at ({62+(\i-1)*22},2)
      {\iflangja{メモリ}{Memory} \i\\\texttt{mem[0x100] = \v}};
    \draw[<->] (q\i.south) -- (dm\i.north);
  }
  % messages to P2
  \draw[->, ln-accent, thick] (dm1.north east) to[bend left=15] (q2.west);
  \draw[->, ln-accent, thick] (dm3.north west) to[bend right=15] (q2.east);
  \node[draw=ln-accent, fill=white, inner sep=1pt, font=\scriptsize\sffamily, text=ln-accent] at (73.5,10.5) {$\boxtimes$};
  \node[draw=ln-accent, fill=white, inner sep=1pt, font=\scriptsize\sffamily, text=ln-accent] at (94.5,10.5) {$\boxtimes$};
  \node[above=1pt of q2, text=ln-accent] {$3+5+8=16$};
  \node[anchor=west, text=ln-muted] at (54,-6) {$\boxtimes$ = \iflangja{メッセージ}{message}};
\end{tikzpicture}
   (body width, 120 mm)
\begin{tikzpicture}[x=1mm, y=1mm, font=\scriptsize\sffamily,
                    >={Stealth[length=1.4mm]},
                    proc/.style={draw=ln-accent, thick, minimum width=13mm, minimum height=6mm},
                    mem/.style={draw=ln-muted, fill=ln-box, minimum height=7mm, align=center}]
  % ---- shared memory
  \node[anchor=west, font=\scriptsize\sffamily\bfseries] at (0,24)
    {\iflangja{共有メモリ}{Shared memory}};
  \foreach \i in {1,2,3} {
    \node[proc] (p\i) at ({8+(\i-1)*16},16) {P\i};
  }
  \node[mem, minimum width=42mm] (sm) at (24,2)
    {\iflangja{メモリ}{Memory}\\\texttt{mem[0x100] = 10}};
  \draw[->] (p1.south) -- (sm.north -| p1.south)
    node[midway, left, align=right, text=ln-muted] {\iflangja{10を\\書く}{writes 10}};
  \draw[<-] (p2.south) -- (sm.north -| p2.south)
    node[midway, right, text=ln-muted] {\iflangja{10を読む}{reads 10}};
  \draw[<->] (p3.south) -- (sm.north -| p3.south);
  % ---- distributed memory
  \node[anchor=west, font=\scriptsize\sffamily\bfseries] at (54,24)
    {\iflangja{分散メモリ}{Distributed memory}};
  \foreach \i/\v/\c in {1/3/orange!45, 2/5/blue!25, 3/8/red!35} {
    \node[proc] (q\i) at ({62+(\i-1)*22},16) {P\i};
    \node[mem, fill=\c, minimum width=17mm] (dm\i) at ({62+(\i-1)*22},2)
      {\iflangja{メモリ}{Memory} \i\\\texttt{mem[0x100] = \v}};
    \draw[<->] (q\i.south) -- (dm\i.north);
  }
  % messages to P2
  \draw[->, ln-accent, thick] (dm1.north east) to[bend left=15] (q2.west);
  \draw[->, ln-accent, thick] (dm3.north west) to[bend right=15] (q2.east);
  \node[draw=ln-accent, fill=white, inner sep=1pt, font=\scriptsize\sffamily, text=ln-accent] at (73.5,10.5) {$\boxtimes$};
  \node[draw=ln-accent, fill=white, inner sep=1pt, font=\scriptsize\sffamily, text=ln-accent] at (94.5,10.5) {$\boxtimes$};
  \node[above=1pt of q2, text=ln-accent] {$3+5+8=16$};
  \node[anchor=west, text=ln-muted] at (54,-6) {$\boxtimes$ = \iflangja{メッセージ}{message}};
\end{tikzpicture}

  \caption{Shared memory (left) and distributed memory (right).  The same
    address \texttt{0x100} names one location on the left and three
    unrelated locations on the right.}
  \label{fig:l02-memory-models}
\end{figure}

\caution{\enquote{Shared memory} in this lesson is the CPU sense, one
address space for all threads.  CUDA's \texttt{\_\_shared\_\_} memory of
Lesson~3 is a small on-chip scratchpad shared only by the threads of one
block.}

Shared memory is the easier model to program, because data sharing needs no
explicit transfer, but it does not scale: every processor contends for the
same memory, and the programmer is responsible for synchronising the
accesses.  Distributed memory scales with the number of processors and
keeps local accesses fast, at the price of making the programmer specify
every transfer.\margin{The largest machines are hybrids: shared memory
inside a node, messages between nodes; the common pairings are MPI with
OpenMP, or MPI with CUDA.  \cite{barney2026parallel}.}

\subsection{OpenMP}

\begin{definition}[OpenMP]
  \term{OpenMP}[OpenMP] is an open standard for shared-memory parallel
  programming in C, C++ and Fortran.  It is a collection of
  \term{compiler directive}[コンパイラ指示文]s (\texttt{\#pragma omp \ldots}),
  library routines (\texttt{omp\_\ldots()}) and
  \term{environment variable}[環境変数]s (\texttt{OMP\_\ldots}) with which
  a serial program is annotated, rather than rewritten, to run in
  parallel.\margin{OpenMP 1.0 appeared in 1997 (Fortran) and 1998 (C/C++);
  the specification cited here is 5.2 (2021), the current one 6.0 (2024).
  \texttt{openmp.org}.}
\end{definition}

Whatever the tool, a shared-memory program has to settle five things: the
\emph{parallelisation} itself (which loop or region runs in parallel), the
\emph{number of threads}, the \emph{scheduling} of work onto those threads,
the \emph{data sharing} policy (which variables are shared, which are
private), and the \emph{synchronisation} that keeps shared data
consistent.  OpenMP has a directive or clause for each; this section
covers the first and \cref{sec:l02-scheduling} the other four.

\subsection{Threads}

\begin{definition}[Thread]
  A \term{thread}[スレッド] is an entity that runs a sequence of
  instructions.  It has its own registers and its own stack, and it shares
  the rest of the address space (code, global data, heap) with the other
  threads of the same process.\margin{On Linux a new thread's stack
  defaults to the \texttt{ulimit -s} soft limit at program start,
  \SI{8}{\mebi\byte} on most distributions;
  \texttt{pthread\_attr\_setstacksize} overrides it per thread.
  \cite{manpages2026pthreadcreate}.}
\end{definition}

A thread is a software concept, not a core: one core may run several
threads, time-sliced by the operating system or interleaved by the
hardware.  The word therefore means
different things on a CPU and on a GPU: a CPU thread is a heavyweight,
independently scheduled flow of control, whereas a GPU thread is one lane of
a warp that executes in lockstep with 31 others (Lesson~1).  This lesson
assumes, for simplicity, one thread per core.\margin{Multithreading gives
a CPU core two or more hardware thread contexts (Lesson~1); a GPU SM holds
up to 64 resident warps on current data-centre GPUs
(\cite[Table 21]{nvidia2024cudaguide}; Lesson~3), each a hardware-scheduled
group of threads.}

\subsection{Parallelising a sum by hand}

The running example is the sum of a ten-element array.  Serially:
\begin{lstlisting}[language=C++, numbers=none]
#include <iostream>
int main() {
  const int size = 10;
  int data[size];
  for (int i = 0; i < size; ++i) data[i] = i + 1;
  int sum = 0;
  for (int i = 0; i < size; ++i) sum += data[i];
  std::cout << "sum = " << sum << '\n';   // 55
  return 0;
}
\end{lstlisting}
Data decomposition over two cores splits the array in half:
\texttt{data[0..4]} goes to a thread on core~1 and \texttt{data[5..9]} to a
thread on core~2 (\cref{fig:l02-data-decomposition}).  Done by hand, the
parallel version has three steps:
\begin{enumerate}
  \item create two threads;
  \item give each thread its half of the array, on which it computes a
    \emph{partial sum} (15 and 40 for the values $1,\ldots,10$);
  \item merge the two partial sums into the total, $15+40=55$.
\end{enumerate}
The merge is a \term*{reduction} (Lesson~1): many values combined by an
associative operation into one.\margin{With thousands of GPU threads the
reduction becomes a tree: partial sums are combined in pairs over
$\log_2$ rounds, within a warp through the shuffle operations of
Lesson~5.}

\begin{figure}[tb]
  \centering
  % Data decomposition of a 10-element sum over two cores, one thread per
% core: each thread adds its half (a partial sum), then the two partial
% sums are reduced to the total.  Redraws lowyx m2l2_data_decomposition2
% and m2l2_manual_parallel.
% Usage: % Data decomposition of a 10-element sum over two cores, one thread per
% core: each thread adds its half (a partial sum), then the two partial
% sums are reduced to the total.  Redraws lowyx m2l2_data_decomposition2
% and m2l2_manual_parallel.
% Usage: % Data decomposition of a 10-element sum over two cores, one thread per
% core: each thread adds its half (a partial sum), then the two partial
% sums are reduced to the total.  Redraws lowyx m2l2_data_decomposition2
% and m2l2_manual_parallel.
% Usage: \input{figures/l02-data-decomposition}   (body width, 120 mm)
\begin{tikzpicture}[x=1mm, y=1mm, font=\scriptsize\sffamily,
                    >={Stealth[length=1.4mm]}]
  \def\cw{9}
  % the array
  \foreach \i in {0,...,9} {
    \pgfmathsetmacro{\v}{int(\i+1)}
    \ifnum\i<5 \def\c{orange!40}\else \def\c{blue!25}\fi
    \fill[\c, draw=ln-muted] (20+\i*\cw,20) rectangle (20+\i*\cw+\cw,26);
    \node at (20+\i*\cw+\cw/2,23) {\v};
    \node[above, text=ln-muted, font=\scriptsize\sffamily] at (20+\i*\cw+\cw/2,26) {data[\i]};
  }
  % braces to cores
  \draw[->, thick] (42.5,19) -- (42.5,13);
  \draw[->, thick] (87.5,19) -- (87.5,13);
  \node[draw=ln-accent, thick, minimum width=30mm, minimum height=9mm, align=center]
    (c1) at (42.5,7) {\iflangja{コア1，スレッド1}{Core 1, thread 1}\\$1+2+3+4+5$};
  \node[draw=ln-accent, thick, minimum width=30mm, minimum height=9mm, align=center]
    (c2) at (87.5,7) {\iflangja{コア2，スレッド2}{Core 2, thread 2}\\$6+7+8+9+10$};
  \node[below=4mm of c1, fill=orange!40, inner sep=2pt] (s1) {\iflangja{部分和}{partial sum} $=15$};
  \node[below=4mm of c2, fill=blue!25, inner sep=2pt] (s2) {\iflangja{部分和}{partial sum} $=40$};
  \draw[->] (c1) -- (s1);
  \draw[->] (c2) -- (s2);
  \node[draw=ln-accent, thick, fill=ln-box, inner sep=3pt] (t) at (65,-16)
    {\iflangja{総和}{total} $=15+40=55$};
  \draw[->] (s1.south) -- (t.west);
  \draw[->] (s2.south) -- (t.east);
  \node[text=ln-muted, anchor=west] at (90,-16) {\iflangja{リダクション}{reduction}};
\end{tikzpicture}
   (body width, 120 mm)
\begin{tikzpicture}[x=1mm, y=1mm, font=\scriptsize\sffamily,
                    >={Stealth[length=1.4mm]}]
  \def\cw{9}
  % the array
  \foreach \i in {0,...,9} {
    \pgfmathsetmacro{\v}{int(\i+1)}
    \ifnum\i<5 \def\c{orange!40}\else \def\c{blue!25}\fi
    \fill[\c, draw=ln-muted] (20+\i*\cw,20) rectangle (20+\i*\cw+\cw,26);
    \node at (20+\i*\cw+\cw/2,23) {\v};
    \node[above, text=ln-muted, font=\scriptsize\sffamily] at (20+\i*\cw+\cw/2,26) {data[\i]};
  }
  % braces to cores
  \draw[->, thick] (42.5,19) -- (42.5,13);
  \draw[->, thick] (87.5,19) -- (87.5,13);
  \node[draw=ln-accent, thick, minimum width=30mm, minimum height=9mm, align=center]
    (c1) at (42.5,7) {\iflangja{コア1，スレッド1}{Core 1, thread 1}\\$1+2+3+4+5$};
  \node[draw=ln-accent, thick, minimum width=30mm, minimum height=9mm, align=center]
    (c2) at (87.5,7) {\iflangja{コア2，スレッド2}{Core 2, thread 2}\\$6+7+8+9+10$};
  \node[below=4mm of c1, fill=orange!40, inner sep=2pt] (s1) {\iflangja{部分和}{partial sum} $=15$};
  \node[below=4mm of c2, fill=blue!25, inner sep=2pt] (s2) {\iflangja{部分和}{partial sum} $=40$};
  \draw[->] (c1) -- (s1);
  \draw[->] (c2) -- (s2);
  \node[draw=ln-accent, thick, fill=ln-box, inner sep=3pt] (t) at (65,-16)
    {\iflangja{総和}{total} $=15+40=55$};
  \draw[->] (s1.south) -- (t.west);
  \draw[->] (s2.south) -- (t.east);
  \node[text=ln-muted, anchor=west] at (90,-16) {\iflangja{リダクション}{reduction}};
\end{tikzpicture}
   (body width, 120 mm)
\begin{tikzpicture}[x=1mm, y=1mm, font=\scriptsize\sffamily,
                    >={Stealth[length=1.4mm]}]
  \def\cw{9}
  % the array
  \foreach \i in {0,...,9} {
    \pgfmathsetmacro{\v}{int(\i+1)}
    \ifnum\i<5 \def\c{orange!40}\else \def\c{blue!25}\fi
    \fill[\c, draw=ln-muted] (20+\i*\cw,20) rectangle (20+\i*\cw+\cw,26);
    \node at (20+\i*\cw+\cw/2,23) {\v};
    \node[above, text=ln-muted, font=\scriptsize\sffamily] at (20+\i*\cw+\cw/2,26) {data[\i]};
  }
  % braces to cores
  \draw[->, thick] (42.5,19) -- (42.5,13);
  \draw[->, thick] (87.5,19) -- (87.5,13);
  \node[draw=ln-accent, thick, minimum width=30mm, minimum height=9mm, align=center]
    (c1) at (42.5,7) {\iflangja{コア1，スレッド1}{Core 1, thread 1}\\$1+2+3+4+5$};
  \node[draw=ln-accent, thick, minimum width=30mm, minimum height=9mm, align=center]
    (c2) at (87.5,7) {\iflangja{コア2，スレッド2}{Core 2, thread 2}\\$6+7+8+9+10$};
  \node[below=4mm of c1, fill=orange!40, inner sep=2pt] (s1) {\iflangja{部分和}{partial sum} $=15$};
  \node[below=4mm of c2, fill=blue!25, inner sep=2pt] (s2) {\iflangja{部分和}{partial sum} $=40$};
  \draw[->] (c1) -- (s1);
  \draw[->] (c2) -- (s2);
  \node[draw=ln-accent, thick, fill=ln-box, inner sep=3pt] (t) at (65,-16)
    {\iflangja{総和}{total} $=15+40=55$};
  \draw[->] (s1.south) -- (t.west);
  \draw[->] (s2.south) -- (t.east);
  \node[text=ln-muted, anchor=west] at (90,-16) {\iflangja{リダクション}{reduction}};
\end{tikzpicture}

  \caption{Data decomposition of the sum: each thread adds its half, then
    the two partial sums are reduced to the total.}
  \label{fig:l02-data-decomposition}
\end{figure}

\subsection{Reduction in shared memory and the race}

In shared memory the reduction is naturally written as each thread adding
its partial sum into one shared variable \texttt{sum}, initially~0.  The
statement \texttt{sum += partial} is three machine operations: load
\texttt{sum}, add, store \texttt{sum}.  If the two threads interleave as in
\cref{tab:l02-race}, both load the initial~0, and the second store
overwrites the first: the result is 40 (or 15, if the stores happen in the
other order) instead of~55.

\begin{table}[tb]
  \centering
  \small
  \caption{An interleaving of two unprotected updates of the shared
    \texttt{sum}; one partial sum is lost.}
  \label{tab:l02-race}
  \begin{tabular}{llll}
    \toprule
    Step & Thread 1 (partial 15) & Thread 2 (partial 40) & \texttt{sum} \\
    \midrule
    1 & load \texttt{sum} $\to 0$ &                            & 0 \\
    2 &                            & load \texttt{sum} $\to 0$ & 0 \\
    3 & add: $0+15=15$             &                            & 0 \\
    4 &                            & add: $0+40=40$             & 0 \\
    5 & store 15                   &                            & 15 \\
    6 &                            & store 40                   & 40 \\
    \bottomrule
  \end{tabular}
\end{table}

\begin{definition}[Data race, critical section]
  A \term{data race}[データ競合] occurs when two threads access the same
  variable without ordering and at least one access is a write; the
  outcome depends on the interleaving.  A piece of code that accesses
  shared data and must therefore be executed by one thread at a time is a
  \term{critical section}[クリティカルセクション].
\end{definition}

\caution{In C and C++ a data race is undefined behaviour, not merely a
wrong number: the compiler may assume it does not happen and, for
instance, keep \texttt{sum} in a register across the whole loop.}

\begin{definition}[Mutex]
  A \term{mutex}[ミューテックス] (from \term{mutual exclusion}[相互排除])
  is a \term{lock}[ロック] that at most one thread can hold.  A thread
  \emph{locks} (acquires) the mutex before entering the critical section
  and \emph{unlocks} (releases) it on leaving; a thread that finds the
  mutex held waits.  With the update of \texttt{sum} inside a mutex, the
  two additions happen one after the other and the result is~55.
\end{definition}

The low-level way to write this on a POSIX system is
\term{POSIX threads (pthreads)}[POSIXスレッド]\index{pthreads|see{POSIX threads (pthreads)}}, a C library of
calls:\margin{Pthreads is IEEE POSIX 1003.1c (1995), a library with no
compiler support; every thread and every lock is explicit.
\cite{barney2026parallel}.}
\begin{itemize}
  \item \texttt{pthread\_create} starts a thread on a given function;
  \item \texttt{pthread\_join} waits for a thread to finish;
  \item \texttt{pthread\_mutex\_lock} and \texttt{pthread\_mutex\_unlock}
    bracket the critical section.
\end{itemize}
Thread creation, data distribution, locking and joining are all the
programmer's code, and each is a place for a bug.

\subsection{The OpenMP version}

OpenMP hides all of that behind one directive.\tip{Compiled without
OpenMP support (no \texttt{-fopenmp}), the pragma is ignored and the
program is the serial one; OpenMP can be added to a working serial program
one loop at a time.}
\begin{lstlisting}[language=C++, numbers=none]
#include <iostream>
#include <omp.h>
int main() {
  const int size = 1000;
  int data[size];
  for (int i = 0; i < size; ++i) data[i] = i + 1;
  int sum = 0;
  #pragma omp parallel for reduction(+:sum)
  for (int i = 0; i < size; ++i) sum += data[i];
  std::cout << "sum = " << sum << '\n';   // 500500
  return 0;
}
\end{lstlisting}
\texttt{parallel for} invokes the loop-parallelism pattern: it creates a
team of threads, divides the iterations among them and joins the threads at
the end of the loop.  \texttt{reduction(+:sum)} tells the compiler that \texttt{sum} is a
reduction over \texttt{+}: every thread gets a private copy initialised to
the identity~0, accumulates into it, and the copies are combined into the
shared \texttt{sum} when the loop ends, with the combination protected
(\texttt{*}, \texttt{max}, \texttt{min} and the bitwise and logical
operators are allowed as well).

\begin{definition}[Directive and clause]
  An OpenMP directive has the form
\begin{lstlisting}[language=C++, numbers=none]
#pragma omp directive [clause [clause] ...]
\end{lstlisting}
  The \textit{directive} names the construct (\texttt{parallel},
  \texttt{for}, \texttt{critical}, \ldots); a \term{clause}[指示節]
  modifies it (\texttt{reduction(+:sum)}, \texttt{num\_threads(4)},
  \texttt{schedule(dynamic)}).  The compiler replaces each directive with
  calls into the OpenMP runtime library, which creates, schedules and joins
  the threads.
\end{definition}

\section{Scheduling, data sharing and synchronisation}
\label{sec:l02-scheduling}

The remaining four tasks of a shared-memory program, the thread count, the
schedule, the sharing policy and the synchronisation, each have an OpenMP
construct.\source{M2, OpenMP vs.\ MPI (part 2); PPP App.~A;
\cite{openmp2021spec}.}

\subsection{How many threads}

The number of threads in a parallel region is set in one of three places:
\begin{itemize}
  \item outside the program, by the environment variable
    \texttt{OMP\_NUM\_THREADS};
  \item inside the program, by the library call
    \texttt{omp\_set\_num\_threads()};
  \item on one region, by the clause \texttt{num\_threads(}$n$\texttt{)}.
\end{itemize}
If nothing is set, the default is the hardware
parallelism of the machine, i.e.\ the number of logical processors
(hardware threads) the operating system reports.\margin{The default is
implementation-defined \cite{openmp2021spec}; GCC's libgomp starts one
thread per CPU available to the process \cite{gcc2026libgomp}, so a
16-core CPU with two hardware threads per core starts 32 threads.}  More threads than cores is allowed: the operating
system time-slices them, which is the situation the next example starts
from.

\subsection{Loop scheduling}

A \texttt{parallel for} must decide which thread executes which
iterations.  The decision matters when the threads do not progress at the
same rate.

\begin{example}[Threads that progress unevenly]
  \label{ex:l02-schedule}
  A loop of \num{600000} iterations of equal cost runs with 6 threads on a
  machine with 4 cores.  The operating system places two threads on each
  of cores 0 and~1 and one thread on each of cores 2 and~3, so the four
  threads on the doubled-up cores each receive half a core.  Let $t$ be the
  time a whole core needs for \num{100000} iterations.
  \begin{itemize}
    \item \emph{Static split}, \num{100000} iterations per thread: the
      threads on cores 2 and~3 finish at time $t$, the four others at
      $2t$.  The loop ends at $2t$, with two cores idle for half of it.
    \item \emph{Ideal}: \num{600000} iterations spread over 4 cores is
      \num{150000} per core, i.e.\ $1.5t$.
    \item \emph{Dynamic, chunks of \num{1000}}: each thread takes the next
      chunk when it finishes one.  A thread with a whole core takes chunks
      twice as often as a thread with half a core, so all six finish within
      about one chunk of $1.5t$, after \num{600} chunk requests.
  \end{itemize}
  The static schedule takes \SI{33}{\percent} longer than the dynamic
  one here, purely from the uneven placement of the threads.
\end{example}

OpenMP's \texttt{schedule(}\textit{kind}\texttt{, }\textit{chunk}\texttt{)}
clause selects the kind of schedule; the three used in practice are
(\cref{fig:l02-schedule}):\margin{Two further kinds: \texttt{auto} leaves
the choice to the compiler or runtime, \texttt{runtime} takes it from the
environment variable \texttt{OMP\_SCHEDULE}.  \cite{openmp2021spec}.}
\begin{description}
  \item[static]  \term{static scheduling}[静的スケジューリング] divides the
    iterations into chunks before the loop starts and assigns them to the
    threads in a fixed round-robin order; with no chunk size, each thread
    gets one contiguous block of about $N/\text{threads}$ iterations.  No
    run-time cost, no adaptation.
  \item[dynamic]  \term{dynamic scheduling}[動的スケジューリング] hands out
    chunks of the given \term{chunk size}[チャンクサイズ] (default~1) on
    request: a thread that finishes its chunk asks for the next.  Chunk
    size~1 adapts best but pays the cost of a request, an atomic update
    of a shared counter, for every iteration; \num{1000} amortises that
    cost over a thousand iterations.
  \item[guided]  \term{guided scheduling}[ガイド付きスケジューリング]
    starts with large chunks and shrinks them as the loop proceeds (for
    instance from \num{1000} down to \num{100}): the early large chunks
    keep the request cost low, the late small ones even out the finish
    times.\margin{The guided chunk is proportional to the remaining
    iterations divided by the number of threads, never below the
    \textit{chunk} argument except for the final chunk.
    \cite{openmp2021spec}.}
\end{description}

\begin{figure}[tb]
  \centering
  % How the three OpenMP loop schedules hand out an iteration space to four
% threads: static gives each thread one contiguous block; dynamic hands out
% fixed chunks on request, so the assignment order is not predictable;
% guided starts with large chunks and shrinks them.  Original.
% Usage: % How the three OpenMP loop schedules hand out an iteration space to four
% threads: static gives each thread one contiguous block; dynamic hands out
% fixed chunks on request, so the assignment order is not predictable;
% guided starts with large chunks and shrinks them.  Original.
% Usage: % How the three OpenMP loop schedules hand out an iteration space to four
% threads: static gives each thread one contiguous block; dynamic hands out
% fixed chunks on request, so the assignment order is not predictable;
% guided starts with large chunks and shrinks them.  Original.
% Usage: \input{figures/l02-schedule}   (body width, 120 mm)
\begin{tikzpicture}[x=1mm, y=1mm, font=\scriptsize\sffamily,
                    >={Stealth[length=1.4mm]}]
  \def\W{88}   % width of the iteration space
  \def\H{5}
  % iteration axis
  \draw[->] (26,13) -- (26+\W+3,13);
  \node[above, text=ln-muted] at (26+\W/2,13) {\iflangja{反復}{iteration}};
  \node[above, text=ln-muted] at (26,13) {0};
  \node[above, text=ln-muted] at (26+\W,13) {$N$};
  % static
  \node[anchor=east] at (24,8) {\texttt{static}};
  \foreach \i/\c in {0/blue!35, 1/orange!55, 2/green!45, 3/red!35} {
    \fill[\c, draw=ln-muted] (26+\i*\W/4,5.5) rectangle (26+\i*\W/4+\W/4,10.5);
    \node at (26+\i*\W/4+\W/8,8) {T\i};
  }
  % dynamic, chunk c: 16 chunks in an irregular thread order
  \node[anchor=east] at (24,0) {\texttt{dynamic,}$c$};
  \foreach \i/\c in {0/blue!35, 1/orange!55, 2/green!45, 3/red!35, 4/orange!55, 5/blue!35,
                     6/red!35, 7/orange!55, 8/green!45, 9/blue!35, 10/orange!55, 11/red!35,
                     12/blue!35, 13/green!45, 14/orange!55, 15/blue!35} {
    \fill[\c, draw=ln-muted] (26+\i*\W/16,-2.5) rectangle (26+\i*\W/16+\W/16,2.5);
  }
  \draw[<->] (26,-4) -- (26+\W/16,-4) node[midway, below, text=ln-muted] {$c$};
  % guided: shrinking chunks
  \node[anchor=east] at (24,-10) {\texttt{guided}};
  \foreach \a/\b/\c in {0/22/blue!35, 22/39/orange!55, 39/52/green!45, 52/62/red!35,
                        62/70/orange!55, 70/76/blue!35, 76/80/green!45, 80/83/red!35,
                        83/85.5/orange!55, 85.5/88/blue!35} {
    \fill[\c, draw=ln-muted] (26+\a,-12.5) rectangle (26+\b,-7.5);
  }
  \node[anchor=west, text=ln-muted, align=left] at (26,-17)
    {\iflangja{色＝そのチャンクを実行するスレッド T0--T3}{colour = the thread T0--T3 that runs the chunk}};
\end{tikzpicture}
   (body width, 120 mm)
\begin{tikzpicture}[x=1mm, y=1mm, font=\scriptsize\sffamily,
                    >={Stealth[length=1.4mm]}]
  \def\W{88}   % width of the iteration space
  \def\H{5}
  % iteration axis
  \draw[->] (26,13) -- (26+\W+3,13);
  \node[above, text=ln-muted] at (26+\W/2,13) {\iflangja{反復}{iteration}};
  \node[above, text=ln-muted] at (26,13) {0};
  \node[above, text=ln-muted] at (26+\W,13) {$N$};
  % static
  \node[anchor=east] at (24,8) {\texttt{static}};
  \foreach \i/\c in {0/blue!35, 1/orange!55, 2/green!45, 3/red!35} {
    \fill[\c, draw=ln-muted] (26+\i*\W/4,5.5) rectangle (26+\i*\W/4+\W/4,10.5);
    \node at (26+\i*\W/4+\W/8,8) {T\i};
  }
  % dynamic, chunk c: 16 chunks in an irregular thread order
  \node[anchor=east] at (24,0) {\texttt{dynamic,}$c$};
  \foreach \i/\c in {0/blue!35, 1/orange!55, 2/green!45, 3/red!35, 4/orange!55, 5/blue!35,
                     6/red!35, 7/orange!55, 8/green!45, 9/blue!35, 10/orange!55, 11/red!35,
                     12/blue!35, 13/green!45, 14/orange!55, 15/blue!35} {
    \fill[\c, draw=ln-muted] (26+\i*\W/16,-2.5) rectangle (26+\i*\W/16+\W/16,2.5);
  }
  \draw[<->] (26,-4) -- (26+\W/16,-4) node[midway, below, text=ln-muted] {$c$};
  % guided: shrinking chunks
  \node[anchor=east] at (24,-10) {\texttt{guided}};
  \foreach \a/\b/\c in {0/22/blue!35, 22/39/orange!55, 39/52/green!45, 52/62/red!35,
                        62/70/orange!55, 70/76/blue!35, 76/80/green!45, 80/83/red!35,
                        83/85.5/orange!55, 85.5/88/blue!35} {
    \fill[\c, draw=ln-muted] (26+\a,-12.5) rectangle (26+\b,-7.5);
  }
  \node[anchor=west, text=ln-muted, align=left] at (26,-17)
    {\iflangja{色＝そのチャンクを実行するスレッド T0--T3}{colour = the thread T0--T3 that runs the chunk}};
\end{tikzpicture}
   (body width, 120 mm)
\begin{tikzpicture}[x=1mm, y=1mm, font=\scriptsize\sffamily,
                    >={Stealth[length=1.4mm]}]
  \def\W{88}   % width of the iteration space
  \def\H{5}
  % iteration axis
  \draw[->] (26,13) -- (26+\W+3,13);
  \node[above, text=ln-muted] at (26+\W/2,13) {\iflangja{反復}{iteration}};
  \node[above, text=ln-muted] at (26,13) {0};
  \node[above, text=ln-muted] at (26+\W,13) {$N$};
  % static
  \node[anchor=east] at (24,8) {\texttt{static}};
  \foreach \i/\c in {0/blue!35, 1/orange!55, 2/green!45, 3/red!35} {
    \fill[\c, draw=ln-muted] (26+\i*\W/4,5.5) rectangle (26+\i*\W/4+\W/4,10.5);
    \node at (26+\i*\W/4+\W/8,8) {T\i};
  }
  % dynamic, chunk c: 16 chunks in an irregular thread order
  \node[anchor=east] at (24,0) {\texttt{dynamic,}$c$};
  \foreach \i/\c in {0/blue!35, 1/orange!55, 2/green!45, 3/red!35, 4/orange!55, 5/blue!35,
                     6/red!35, 7/orange!55, 8/green!45, 9/blue!35, 10/orange!55, 11/red!35,
                     12/blue!35, 13/green!45, 14/orange!55, 15/blue!35} {
    \fill[\c, draw=ln-muted] (26+\i*\W/16,-2.5) rectangle (26+\i*\W/16+\W/16,2.5);
  }
  \draw[<->] (26,-4) -- (26+\W/16,-4) node[midway, below, text=ln-muted] {$c$};
  % guided: shrinking chunks
  \node[anchor=east] at (24,-10) {\texttt{guided}};
  \foreach \a/\b/\c in {0/22/blue!35, 22/39/orange!55, 39/52/green!45, 52/62/red!35,
                        62/70/orange!55, 70/76/blue!35, 76/80/green!45, 80/83/red!35,
                        83/85.5/orange!55, 85.5/88/blue!35} {
    \fill[\c, draw=ln-muted] (26+\a,-12.5) rectangle (26+\b,-7.5);
  }
  \node[anchor=west, text=ln-muted, align=left] at (26,-17)
    {\iflangja{色＝そのチャンクを実行するスレッド T0--T3}{colour = the thread T0--T3 that runs the chunk}};
\end{tikzpicture}

  \caption{Three loop schedules for four threads.  Static assigns one
    block per thread in advance; dynamic hands out fixed chunks in the
    order the threads ask for them; guided hands out shrinking chunks.}
  \label{fig:l02-schedule}
\end{figure}

Dynamic and guided scheduling adapt to anything that makes threads
unequal at run time: uneven placement as in \cref{ex:l02-schedule}, a
slower machine in a heterogeneous cluster, iterations of varying cost, or
interference from other programs.  Keeping all threads busy until the end
is \term{load balancing}[負荷分散]; the time of a parallel loop is the
time of its slowest thread.\tip{Regular loops with equal iterations:
\texttt{static}.  Irregular or unknown cost: \texttt{dynamic} with a chunk
large enough that a request costs far less than a chunk; \texttt{guided}
when the iteration count is large.}

\subsection{Data sharing}

In the sum, each thread's partial sum must be its own copy, while the
total is one location.  OpenMP calls the first a
\term{private variable}[プライベート変数] (one instance per thread, as
with the \texttt{reduction} copies) and the second a
\term{shared variable}[共有変数] (one instance seen by all).  The
programmer states the policy with the \texttt{private(\ldots)} and
\texttt{shared(\ldots)} clauses; by default, variables declared outside the
parallel region are shared and the loop index is
private.\caution{A variable left shared by the default and written in the
loop is a data race.  \texttt{default(none)} forces every variable to be
listed explicitly and turns the omission into a compile error.}

\subsection{Synchronisation}

Synchronisation is what makes the shared-memory program correct: it orders
the accesses of different threads to shared data.  OpenMP provides three
constructs.

\begin{description}
  \item[Barrier]  A \term{barrier}[バリア] is a point that every thread
    of the team must reach before any of them continues
    (\cref{fig:l02-barrier}): the work after the barrier (green) does not
    start until all the work before it (blue) has finished.  It is written
    \texttt{\#pragma omp barrier}; the end of a \texttt{parallel} region
    is always an implicit barrier, and so is the end of a \texttt{for},
    \texttt{sections} or \texttt{single} construct unless a
    \texttt{nowait} clause removes it.\margin{CUDA's
    \texttt{\_\_syncthreads()} (Lesson~3) is a barrier over the threads
    of one block; a barrier over the whole grid exists only through
    cooperative groups with a cooperative launch (Lesson~5).}
  \item[Critical]  \texttt{\#pragma omp critical [(\textit{name})]} makes
    the following block a critical section: one thread at a time executes
    it, the others wait at its entry (\cref{fig:l02-critical}).  With four
    threads in the region, the four executions of the block are
    serialised back to back.  Critical sections with different names may
    run concurrently; unnamed ones all share one lock.
\begin{lstlisting}[language=C++, numbers=none]
#pragma omp parallel num_threads(4)
{
  // parallel work ...
  #pragma omp critical
  {
    shared_total += local_value;   // one thread at a time
  }
}
\end{lstlisting}
  \item[Atomic]  \texttt{\#pragma omp atomic} makes the single update
    that follows it an \term{atomic operation}[アトミック操作]: its load,
    add and store happen as a whole or not at all, so no other thread's
    update can slip between them.  It is implemented with a hardware
    read-modify-write instruction where one exists and with a lock
    otherwise, and is cheaper than a critical section when it
    applies.\margin{CUDA's \texttt{atomicAdd()} plays the same role on
    the GPU; Lesson~5 lists it among the operations to use sparingly.}
\begin{lstlisting}[language=C++, numbers=none]
#pragma omp parallel for
for (int i = 0; i < n; ++i) {
  #pragma omp atomic
  counter++;          // load, add, store as one unit
}
\end{lstlisting}
\end{description}

\begin{figure}[tb]
  \centering
  % A barrier: five threads finish the first phase (blue) at different
% times; none starts the second phase (green) before all have arrived.
% Redraws lowyx m2l3_barrier.
% Usage: % A barrier: five threads finish the first phase (blue) at different
% times; none starts the second phase (green) before all have arrived.
% Redraws lowyx m2l3_barrier.
% Usage: % A barrier: five threads finish the first phase (blue) at different
% times; none starts the second phase (green) before all have arrived.
% Redraws lowyx m2l3_barrier.
% Usage: \input{figures/l02-barrier}   (margin width, 50 mm)
\begin{tikzpicture}[x=1mm, y=1mm, font=\scriptsize\sffamily,
                    >={Stealth[length=1.4mm]}]
  \fill[blue!15] (0,-1) rectangle (24,14);
  \fill[green!20] (28,-1) rectangle (48,14);
  \fill[yellow!70, draw=ln-muted] (24,-2) rectangle (28,15);
  \node[above, text=ln-muted] at (26,15) {\iflangja{バリア}{barrier}};
  \foreach \y/\e in {12/20, 9/14, 6/23, 3/17, 0/10} {
    \draw[thick, blue!60!black] (0,\y) -- (\e,\y);
    \draw[densely dotted, blue!60!black] (\e,\y) -- (24,\y);
    \draw[->, thick, green!50!black] (28,\y) -- (46,\y);
  }
  \node[below, text=ln-muted] at (12,-2) {\iflangja{フェーズ1（待ち）}{phase 1 (then wait)}};
  \node[below, text=ln-muted] at (38,-2) {\iflangja{フェーズ2}{phase 2}};
  \draw[->] (0,-8) -- (48,-8) node[below left, text=ln-muted] {\iflangja{時間}{time}};
\end{tikzpicture}
   (margin width, 50 mm)
\begin{tikzpicture}[x=1mm, y=1mm, font=\scriptsize\sffamily,
                    >={Stealth[length=1.4mm]}]
  \fill[blue!15] (0,-1) rectangle (24,14);
  \fill[green!20] (28,-1) rectangle (48,14);
  \fill[yellow!70, draw=ln-muted] (24,-2) rectangle (28,15);
  \node[above, text=ln-muted] at (26,15) {\iflangja{バリア}{barrier}};
  \foreach \y/\e in {12/20, 9/14, 6/23, 3/17, 0/10} {
    \draw[thick, blue!60!black] (0,\y) -- (\e,\y);
    \draw[densely dotted, blue!60!black] (\e,\y) -- (24,\y);
    \draw[->, thick, green!50!black] (28,\y) -- (46,\y);
  }
  \node[below, text=ln-muted] at (12,-2) {\iflangja{フェーズ1（待ち）}{phase 1 (then wait)}};
  \node[below, text=ln-muted] at (38,-2) {\iflangja{フェーズ2}{phase 2}};
  \draw[->] (0,-8) -- (48,-8) node[below left, text=ln-muted] {\iflangja{時間}{time}};
\end{tikzpicture}
   (margin width, 50 mm)
\begin{tikzpicture}[x=1mm, y=1mm, font=\scriptsize\sffamily,
                    >={Stealth[length=1.4mm]}]
  \fill[blue!15] (0,-1) rectangle (24,14);
  \fill[green!20] (28,-1) rectangle (48,14);
  \fill[yellow!70, draw=ln-muted] (24,-2) rectangle (28,15);
  \node[above, text=ln-muted] at (26,15) {\iflangja{バリア}{barrier}};
  \foreach \y/\e in {12/20, 9/14, 6/23, 3/17, 0/10} {
    \draw[thick, blue!60!black] (0,\y) -- (\e,\y);
    \draw[densely dotted, blue!60!black] (\e,\y) -- (24,\y);
    \draw[->, thick, green!50!black] (28,\y) -- (46,\y);
  }
  \node[below, text=ln-muted] at (12,-2) {\iflangja{フェーズ1（待ち）}{phase 1 (then wait)}};
  \node[below, text=ln-muted] at (38,-2) {\iflangja{フェーズ2}{phase 2}};
  \draw[->] (0,-8) -- (48,-8) node[below left, text=ln-muted] {\iflangja{時間}{time}};
\end{tikzpicture}

  \caption{A barrier: no thread starts phase 2 before all five have
    finished phase 1.}
  \label{fig:l02-barrier}
\end{figure}

\begin{figure}[tb]
  \centering
  % A critical section with four threads: each thread runs its critical
% section (red) once, and the four executions never overlap in time.
% Redraws lowyx m2l3_critical.
% Usage: % A critical section with four threads: each thread runs its critical
% section (red) once, and the four executions never overlap in time.
% Redraws lowyx m2l3_critical.
% Usage: % A critical section with four threads: each thread runs its critical
% section (red) once, and the four executions never overlap in time.
% Redraws lowyx m2l3_critical.
% Usage: \input{figures/l02-critical}   (body width, 120 mm)
\begin{tikzpicture}[x=1mm, y=1mm, font=\scriptsize\sffamily,
                    >={Stealth[length=1.4mm]}]
  \foreach \i/\s in {1/10, 2/22, 3/34, 4/46} {
    \pgfmathsetmacro{\y}{-\i*6}
    \node[anchor=east] at (-1,\y) {Th \i};
    \fill[blue!25] (0,\y-2) rectangle (70,\y+2);
    \fill[red!60, draw=red!70!black] (\s,\y-2) rectangle (\s+10,\y+2);
    \node[text=white, font=\scriptsize\sffamily\bfseries] at (\s+5,\y) {CS};
  }
  \draw[->] (0,-28) -- (72,-28) node[below left, text=ln-muted] {\iflangja{時間}{time}};
  \fill[blue!25] (78,-10) rectangle (82,-7);
  \node[anchor=west] at (83,-8.5) {\iflangja{並列部分}{parallel work}};
  \fill[red!60] (78,-16) rectangle (82,-13);
  \node[anchor=west] at (83,-14.5) {\iflangja{クリティカルセクション}{critical section}};
\end{tikzpicture}
   (body width, 120 mm)
\begin{tikzpicture}[x=1mm, y=1mm, font=\scriptsize\sffamily,
                    >={Stealth[length=1.4mm]}]
  \foreach \i/\s in {1/10, 2/22, 3/34, 4/46} {
    \pgfmathsetmacro{\y}{-\i*6}
    \node[anchor=east] at (-1,\y) {Th \i};
    \fill[blue!25] (0,\y-2) rectangle (70,\y+2);
    \fill[red!60, draw=red!70!black] (\s,\y-2) rectangle (\s+10,\y+2);
    \node[text=white, font=\scriptsize\sffamily\bfseries] at (\s+5,\y) {CS};
  }
  \draw[->] (0,-28) -- (72,-28) node[below left, text=ln-muted] {\iflangja{時間}{time}};
  \fill[blue!25] (78,-10) rectangle (82,-7);
  \node[anchor=west] at (83,-8.5) {\iflangja{並列部分}{parallel work}};
  \fill[red!60] (78,-16) rectangle (82,-13);
  \node[anchor=west] at (83,-14.5) {\iflangja{クリティカルセクション}{critical section}};
\end{tikzpicture}
   (body width, 120 mm)
\begin{tikzpicture}[x=1mm, y=1mm, font=\scriptsize\sffamily,
                    >={Stealth[length=1.4mm]}]
  \foreach \i/\s in {1/10, 2/22, 3/34, 4/46} {
    \pgfmathsetmacro{\y}{-\i*6}
    \node[anchor=east] at (-1,\y) {Th \i};
    \fill[blue!25] (0,\y-2) rectangle (70,\y+2);
    \fill[red!60, draw=red!70!black] (\s,\y-2) rectangle (\s+10,\y+2);
    \node[text=white, font=\scriptsize\sffamily\bfseries] at (\s+5,\y) {CS};
  }
  \draw[->] (0,-28) -- (72,-28) node[below left, text=ln-muted] {\iflangja{時間}{time}};
  \fill[blue!25] (78,-10) rectangle (82,-7);
  \node[anchor=west] at (83,-8.5) {\iflangja{並列部分}{parallel work}};
  \fill[red!60] (78,-16) rectangle (82,-13);
  \node[anchor=west] at (83,-14.5) {\iflangja{クリティカルセクション}{critical section}};
\end{tikzpicture}

  \caption{Four threads and one critical section: the red executions never
    overlap, the rest of the work does.}
  \label{fig:l02-critical}
\end{figure}

\caution{\texttt{atomic} covers one scalar update of the forms
\texttt{x++}, \texttt{x += e} and the like.  Two statements, or an update
of two variables that must stay consistent, need \texttt{critical}.}

\subsection{Work outside loops}

Not all parallel work is a loop: \term{parallel sections}[並列セクション]
runs independent blocks of code concurrently, one thread per
section:\margin{OpenMP 3.0 (2008) added \texttt{\#pragma omp task}, which
generalises \texttt{sections} to a dynamic set of tasks, the master/worker
pattern of \cref{sec:l02-patterns}.  \cite{openmp2021spec}.}
\begin{lstlisting}[language=C++, numbers=none]
#pragma omp parallel sections
{
  #pragma omp section
  { work_1(); }
  #pragma omp section
  { work_2(); }
}
\end{lstlisting}
Two further constructs restrict, rather than create, parallelism:
\begin{description}
  \item[ordered]  Inside a loop declared \texttt{for ordered}, the block
    marked \texttt{\#pragma omp ordered} executes in iteration order even
    though the rest of each iteration runs in parallel.\caution{Everything
    inside the \texttt{ordered} block is serialised; a loop whose body is
    mostly ordered gains nothing from being parallel.}  A loop that prints
    \enquote{thread $k$ is doing iteration $i$} from five threads prints
    its lines in an arbitrary order (1, 2, 4, 3, 5, say) without
    \texttt{ordered} and as 1, 2, 3, 4, 5 with it.
  \item[single]  \texttt{\#pragma omp single} runs the block on exactly
    one thread of the team, unspecified which, while the others skip it
    and wait at the implicit barrier after it; typical for a one-time
    initialisation or a message inside a parallel region.\tip{When the
    other threads need not wait for the \texttt{single} block,
    \texttt{single nowait} drops its barrier; the next \texttt{for} or
    \texttt{barrier} still orders everything.}
\end{description}

\begin{keypoint}
  OpenMP makes the five decisions of a shared-memory program explicit:
  \texttt{parallel for} (what), \texttt{OMP\_NUM\_THREADS} (how many),
  \texttt{schedule} (who does which iteration), \texttt{private}/\texttt{shared}
  (whose data), and \texttt{barrier}/\texttt{critical}/\texttt{atomic}
  (in what order).  The same five decisions return in CUDA as the kernel,
  the launch configuration, the block scheduler, the memory spaces and
  \texttt{\_\_syncthreads()}.
\end{keypoint}

\section{Programming with MPI}
\label{sec:l02-mpi}

CUDA programming combines the two models of \cref{sec:l02-openmp}: some
memory regions are shared among all the cores of a GPU, while others are
invisible from outside and need explicit communication.  A GPU programmer
therefore needs the distributed-memory habits as much as the shared-memory
ones, which is why message passing is introduced here.\source{M2,
programming with MPI; PPP App.~B; \cite{barney2026parallel}, message
passing model; PMPP ch.~20.}

\subsection{Processes and messages}

\begin{definition}[MPI]
  \term{MPI (Message Passing Interface)}[MPI]\index{Message Passing Interface|see{MPI (Message Passing Interface)}} is the standard library for
  \term{message passing}[メッセージパッシング]: a set of
  \term{process}[プロセス]es, each with its own address space, on the
  same node or on different nodes, exchange data only by sending and
  receiving messages.
  Each process is identified by its \term{rank}[ランク], an integer from
  0 to the number of processes minus~1.\margin{MPI Forum 1992; MPI-1
  1994; version 4.1 (2023) is cited here, 5.0 (2025) is current.
  \cite{barney2026parallel}.}
\end{definition}

Two kinds of operation cover most of the library
(\cref{fig:l02-mpi-messages}):
\begin{description}
  \item[Point-to-point]  \term{point-to-point communication}[1対1通信]
    involves exactly two processes:\margin{A message costs a latency
    (sending a 0-byte message, typically microseconds) plus the size
    divided by the bandwidth; many small messages are far worse than one
    large one.  \cite{barney2026parallel}.} process~0 calls
    \texttt{MPI\_Send()} with an integer and the destination rank~1;
    process~1 calls \texttt{MPI\_Recv()} naming source rank~0 and obtains
    the value.  The message is delivered only to a matching receive posted
    by the destination; both sides take part in every transfer.
  \item[Collective]  \term{collective communication}[集団通信] involves
    every process of a group at once.  The simplest is the
    \term{broadcast}[ブロードキャスト], \texttt{MPI\_Bcast()}: the root
    process's value is delivered to all the other processes; all of them
    call the same function.\margin{Other collectives: \emph{scatter}
    (split one array among the ranks), \emph{gather} (the reverse) and
    \emph{reduce} (combine one value per rank, e.g.\ a global sum).
    \cite{mpi2023standard}.}
\end{description}

\begin{figure}[tb]
  \centering
  % MPI communication: point-to-point (left, MPI_Send on rank 0 matched by
% MPI_Recv on rank 1) and a broadcast (right, MPI_Bcast delivers the root's
% value to every other rank).  Redraws lowyx m2l4_mpi_programming and
% m2l4_broadcasting.
% Usage: % MPI communication: point-to-point (left, MPI_Send on rank 0 matched by
% MPI_Recv on rank 1) and a broadcast (right, MPI_Bcast delivers the root's
% value to every other rank).  Redraws lowyx m2l4_mpi_programming and
% m2l4_broadcasting.
% Usage: % MPI communication: point-to-point (left, MPI_Send on rank 0 matched by
% MPI_Recv on rank 1) and a broadcast (right, MPI_Bcast delivers the root's
% value to every other rank).  Redraws lowyx m2l4_mpi_programming and
% m2l4_broadcasting.
% Usage: \input{figures/l02-mpi-messages}   (body width, 120 mm)
\begin{tikzpicture}[x=1mm, y=1mm, font=\scriptsize\sffamily,
                    >={Stealth[length=1.4mm]},
                    proc/.style={draw=ln-accent, thick, minimum width=17mm, minimum height=7mm, align=center},
                    env/.style={draw=ln-accent, fill=white, inner sep=1pt, text=ln-accent}]
  % point-to-point
  \node[anchor=west, font=\scriptsize\sffamily\bfseries] at (0,16)
    {\iflangja{1対1通信}{Point-to-point}};
  \node[proc] (a) at (9,6) {\iflangja{プロセス}{Process} 0\\\texttt{MPI\_Send}};
  \node[proc] (b) at (43,6) {\iflangja{プロセス}{Process} 1\\\texttt{MPI\_Recv}};
  \draw[->, thick, ln-accent] (a.east) -- (b.west);
  \node[env] at (26,6) {$\boxtimes$};
  \node[below=1pt of a, text=ln-muted] {\iflangja{ランク}{rank} 0};
  \node[below=1pt of b, text=ln-muted] {\iflangja{ランク}{rank} 1};
  % broadcast
  \node[anchor=west, font=\scriptsize\sffamily\bfseries] at (62,16)
    {\iflangja{ブロードキャスト}{Broadcast}};
  \node[proc] (r) at (73,6) {\iflangja{プロセス}{Process} 0\\\texttt{MPI\_Bcast}};
  \foreach \i/\y in {1/14, 2/6, 3/-2} {
    \node[proc, minimum width=15mm] (d\i) at (108,\y) {\iflangja{プロセス}{Process} \i};
    \draw[->, thick, ln-accent] (r.east) -- (d\i.west);
    \node[env] at (92,{6+(\y-6)*0.52}) {$\boxtimes$};
  }
  \node[below=1pt of r, text=ln-muted] {\iflangja{ルート}{root}};
\end{tikzpicture}
   (body width, 120 mm)
\begin{tikzpicture}[x=1mm, y=1mm, font=\scriptsize\sffamily,
                    >={Stealth[length=1.4mm]},
                    proc/.style={draw=ln-accent, thick, minimum width=17mm, minimum height=7mm, align=center},
                    env/.style={draw=ln-accent, fill=white, inner sep=1pt, text=ln-accent}]
  % point-to-point
  \node[anchor=west, font=\scriptsize\sffamily\bfseries] at (0,16)
    {\iflangja{1対1通信}{Point-to-point}};
  \node[proc] (a) at (9,6) {\iflangja{プロセス}{Process} 0\\\texttt{MPI\_Send}};
  \node[proc] (b) at (43,6) {\iflangja{プロセス}{Process} 1\\\texttt{MPI\_Recv}};
  \draw[->, thick, ln-accent] (a.east) -- (b.west);
  \node[env] at (26,6) {$\boxtimes$};
  \node[below=1pt of a, text=ln-muted] {\iflangja{ランク}{rank} 0};
  \node[below=1pt of b, text=ln-muted] {\iflangja{ランク}{rank} 1};
  % broadcast
  \node[anchor=west, font=\scriptsize\sffamily\bfseries] at (62,16)
    {\iflangja{ブロードキャスト}{Broadcast}};
  \node[proc] (r) at (73,6) {\iflangja{プロセス}{Process} 0\\\texttt{MPI\_Bcast}};
  \foreach \i/\y in {1/14, 2/6, 3/-2} {
    \node[proc, minimum width=15mm] (d\i) at (108,\y) {\iflangja{プロセス}{Process} \i};
    \draw[->, thick, ln-accent] (r.east) -- (d\i.west);
    \node[env] at (92,{6+(\y-6)*0.52}) {$\boxtimes$};
  }
  \node[below=1pt of r, text=ln-muted] {\iflangja{ルート}{root}};
\end{tikzpicture}
   (body width, 120 mm)
\begin{tikzpicture}[x=1mm, y=1mm, font=\scriptsize\sffamily,
                    >={Stealth[length=1.4mm]},
                    proc/.style={draw=ln-accent, thick, minimum width=17mm, minimum height=7mm, align=center},
                    env/.style={draw=ln-accent, fill=white, inner sep=1pt, text=ln-accent}]
  % point-to-point
  \node[anchor=west, font=\scriptsize\sffamily\bfseries] at (0,16)
    {\iflangja{1対1通信}{Point-to-point}};
  \node[proc] (a) at (9,6) {\iflangja{プロセス}{Process} 0\\\texttt{MPI\_Send}};
  \node[proc] (b) at (43,6) {\iflangja{プロセス}{Process} 1\\\texttt{MPI\_Recv}};
  \draw[->, thick, ln-accent] (a.east) -- (b.west);
  \node[env] at (26,6) {$\boxtimes$};
  \node[below=1pt of a, text=ln-muted] {\iflangja{ランク}{rank} 0};
  \node[below=1pt of b, text=ln-muted] {\iflangja{ランク}{rank} 1};
  % broadcast
  \node[anchor=west, font=\scriptsize\sffamily\bfseries] at (62,16)
    {\iflangja{ブロードキャスト}{Broadcast}};
  \node[proc] (r) at (73,6) {\iflangja{プロセス}{Process} 0\\\texttt{MPI\_Bcast}};
  \foreach \i/\y in {1/14, 2/6, 3/-2} {
    \node[proc, minimum width=15mm] (d\i) at (108,\y) {\iflangja{プロセス}{Process} \i};
    \draw[->, thick, ln-accent] (r.east) -- (d\i.west);
    \node[env] at (92,{6+(\y-6)*0.52}) {$\boxtimes$};
  }
  \node[below=1pt of r, text=ln-muted] {\iflangja{ルート}{root}};
\end{tikzpicture}

  \caption{Point-to-point (left) and collective (right) communication in
    MPI.}
  \label{fig:l02-mpi-messages}
\end{figure}

A minimal program shows both calls.  The same executable is started
once per process (\texttt{mpirun -np 4 ./a.out} starts four), every copy
runs it from the top, and the rank is the only thing that tells the
copies apart and decides what each one does (SPMD again).
\begin{lstlisting}[language=C, numbers=none]
#include <mpi.h>
int main(int argc, char **argv) {
  int rank, value = 0;
  MPI_Init(&argc, &argv);
  MPI_Comm_rank(MPI_COMM_WORLD, &rank);
  if (rank == 0) {
    value = 42;
    MPI_Send(&value, 1, MPI_INT, 1, 0, MPI_COMM_WORLD);
  } else if (rank == 1) {
    MPI_Recv(&value, 1, MPI_INT, 0, 0, MPI_COMM_WORLD,
             MPI_STATUS_IGNORE);
  }
  MPI_Bcast(&value, 1, MPI_INT, 0, MPI_COMM_WORLD);
  MPI_Finalize();
  return 0;
}
\end{lstlisting}

\subsection{A stencil over distributed data}

\begin{definition}[Stencil]
  A \term{stencil}[ステンシル] computation updates every element of a
  grid from a fixed pattern of its neighbours.  The example here replaces
  each element $C$ by the average of its four neighbours, north, west,
  east and south (\cref{fig:l02-stencil-halo}, left):
  \begin{equation}
    C' = \frac{N + W + E + S}{4}.
    \label{eq:l02-stencil}
  \end{equation}
\end{definition}

Stencils are common in high-performance computing (finite-difference
solvers, image filters).\margin{Stencils are named by their point count:
this is the 5-point stencil (most variants read the centre too), and
\cref{eq:l02-stencil} is one Jacobi sweep for the discrete Laplace
equation, i.e.\ steady-state heat conduction.  Image convolution with a
$3\times3$ filter is a 9-point stencil, Lesson~3.}

When the grid is distributed over four processes, each process owns one
quadrant and can access only that (\cref{fig:l02-stencil-halo}, left).  An
element $C$ at the edge of process~2's region needs neighbour~1, which
belongs to process~0, and neighbour~3, which belongs to process~3.  No
load can fetch them; they must arrive by message.

\begin{figure}[tb]
  \centering
  % A four-point stencil on a 10 x 8 grid split among four processes (left):
% the element C in the blue region needs neighbour 1 from the orange region
% and neighbour 3 from the green region.  Right: the blue region with its
% one-cell halo (grey) received from the neighbouring processes.
% Redraws lowyx m2l4_stencil, m2l4_stencil2 and m2l4_stencil3.
% Usage: % A four-point stencil on a 10 x 8 grid split among four processes (left):
% the element C in the blue region needs neighbour 1 from the orange region
% and neighbour 3 from the green region.  Right: the blue region with its
% one-cell halo (grey) received from the neighbouring processes.
% Redraws lowyx m2l4_stencil, m2l4_stencil2 and m2l4_stencil3.
% Usage: % A four-point stencil on a 10 x 8 grid split among four processes (left):
% the element C in the blue region needs neighbour 1 from the orange region
% and neighbour 3 from the green region.  Right: the blue region with its
% one-cell halo (grey) received from the neighbouring processes.
% Redraws lowyx m2l4_stencil, m2l4_stencil2 and m2l4_stencil3.
% Usage: \input{figures/l02-stencil-halo}   (body width, 120 mm)
\begin{tikzpicture}[x=1mm, y=1mm, font=\scriptsize\sffamily,
                    >={Stealth[length=1.4mm]}]
  \def\s{3.4}  % cell size
  % ---- left: full grid, four quadrants (rows counted from the top)
  \foreach \r in {1,...,8} {
    \foreach \c in {1,...,10} {
      \ifnum\r<5
        \ifnum\c<6 \def\col{orange!45}\else \def\col{yellow!60}\fi
      \else
        \ifnum\c<6 \def\col{blue!30}\else \def\col{green!45}\fi
      \fi
      \fill[\col, draw=white, line width=0.3pt] ({(\c-1)*\s},{-\r*\s}) rectangle ({\c*\s},{-(\r-1)*\s});
    }
  }
  \draw[ln-muted] (0,0) rectangle (10*\s,-8*\s);
  \draw[ln-muted, thick] (5*\s,0) -- (5*\s,-8*\s);
  \draw[ln-muted, thick] (0,-4*\s) -- (10*\s,-4*\s);
  \foreach \c/\r/\t in {5/5/C, 5/4/1, 4/5/2, 6/5/3, 5/6/4} {
    \node[font=\scriptsize\sffamily\bfseries] at ({(\c-0.5)*\s},{-(\r-0.5)*\s}) {\t};
  }
  \node[text=ln-muted] at (2.5*\s,1.8) {P0};
  \node[text=ln-muted] at (7.5*\s,1.8) {P1};
  \node[text=ln-muted] at (-2.5,-6*\s) {P2};
  \node[text=ln-muted] at (10*\s+2.5,-6*\s) {P3};
  \node[anchor=north, text=ln-muted, align=center] at (5*\s,-8*\s-1)
    {\iflangja{4プロセスに分割した格子}{grid split among four processes}};
  % ---- right: the blue region plus its halo
  \begin{scope}[shift={(52,0)}]
    \foreach \c in {1,...,5} {
      \fill[gray!45, draw=white, line width=0.3pt] ({(\c-1)*\s},{-4*\s}) rectangle ({\c*\s},{-3*\s});
    }
    \foreach \r in {5,...,8} {
      \fill[gray!45, draw=white, line width=0.3pt] ({5*\s},{-\r*\s}) rectangle ({6*\s},{-(\r-1)*\s});
      \foreach \c in {1,...,5} {
        \fill[blue!30, draw=white, line width=0.3pt] ({(\c-1)*\s},{-\r*\s}) rectangle ({\c*\s},{-(\r-1)*\s});
      }
    }
    \draw[ln-muted] (0,-4*\s) rectangle (5*\s,-8*\s);
    \foreach \c/\r/\t in {5/5/C, 5/4/1, 4/5/2, 6/5/3, 5/6/4} {
      \node[font=\scriptsize\sffamily\bfseries] at ({(\c-0.5)*\s},{-(\r-0.5)*\s}) {\t};
    }
    \draw[->, thick, ln-accent] (2.5*\s,-1.2*\s) -- (2.5*\s,-2.9*\s)
      node[pos=0, above, text=ln-accent] {\iflangja{P0から}{from P0}};
    \draw[->, thick, ln-accent] (8*\s,-6.5*\s) -- (6.2*\s,-6.5*\s)
      node[pos=0, right, text=ln-accent] {\iflangja{P3から}{from P3}};
    \node[anchor=north, text=ln-muted, align=center] at (3*\s,-8*\s-1)
      {\iflangja{P2の領域とハロー（灰色）}{P2's region and its halo (grey)}};
  \end{scope}
\end{tikzpicture}
   (body width, 120 mm)
\begin{tikzpicture}[x=1mm, y=1mm, font=\scriptsize\sffamily,
                    >={Stealth[length=1.4mm]}]
  \def\s{3.4}  % cell size
  % ---- left: full grid, four quadrants (rows counted from the top)
  \foreach \r in {1,...,8} {
    \foreach \c in {1,...,10} {
      \ifnum\r<5
        \ifnum\c<6 \def\col{orange!45}\else \def\col{yellow!60}\fi
      \else
        \ifnum\c<6 \def\col{blue!30}\else \def\col{green!45}\fi
      \fi
      \fill[\col, draw=white, line width=0.3pt] ({(\c-1)*\s},{-\r*\s}) rectangle ({\c*\s},{-(\r-1)*\s});
    }
  }
  \draw[ln-muted] (0,0) rectangle (10*\s,-8*\s);
  \draw[ln-muted, thick] (5*\s,0) -- (5*\s,-8*\s);
  \draw[ln-muted, thick] (0,-4*\s) -- (10*\s,-4*\s);
  \foreach \c/\r/\t in {5/5/C, 5/4/1, 4/5/2, 6/5/3, 5/6/4} {
    \node[font=\scriptsize\sffamily\bfseries] at ({(\c-0.5)*\s},{-(\r-0.5)*\s}) {\t};
  }
  \node[text=ln-muted] at (2.5*\s,1.8) {P0};
  \node[text=ln-muted] at (7.5*\s,1.8) {P1};
  \node[text=ln-muted] at (-2.5,-6*\s) {P2};
  \node[text=ln-muted] at (10*\s+2.5,-6*\s) {P3};
  \node[anchor=north, text=ln-muted, align=center] at (5*\s,-8*\s-1)
    {\iflangja{4プロセスに分割した格子}{grid split among four processes}};
  % ---- right: the blue region plus its halo
  \begin{scope}[shift={(52,0)}]
    \foreach \c in {1,...,5} {
      \fill[gray!45, draw=white, line width=0.3pt] ({(\c-1)*\s},{-4*\s}) rectangle ({\c*\s},{-3*\s});
    }
    \foreach \r in {5,...,8} {
      \fill[gray!45, draw=white, line width=0.3pt] ({5*\s},{-\r*\s}) rectangle ({6*\s},{-(\r-1)*\s});
      \foreach \c in {1,...,5} {
        \fill[blue!30, draw=white, line width=0.3pt] ({(\c-1)*\s},{-\r*\s}) rectangle ({\c*\s},{-(\r-1)*\s});
      }
    }
    \draw[ln-muted] (0,-4*\s) rectangle (5*\s,-8*\s);
    \foreach \c/\r/\t in {5/5/C, 5/4/1, 4/5/2, 6/5/3, 5/6/4} {
      \node[font=\scriptsize\sffamily\bfseries] at ({(\c-0.5)*\s},{-(\r-0.5)*\s}) {\t};
    }
    \draw[->, thick, ln-accent] (2.5*\s,-1.2*\s) -- (2.5*\s,-2.9*\s)
      node[pos=0, above, text=ln-accent] {\iflangja{P0から}{from P0}};
    \draw[->, thick, ln-accent] (8*\s,-6.5*\s) -- (6.2*\s,-6.5*\s)
      node[pos=0, right, text=ln-accent] {\iflangja{P3から}{from P3}};
    \node[anchor=north, text=ln-muted, align=center] at (3*\s,-8*\s-1)
      {\iflangja{P2の領域とハロー（灰色）}{P2's region and its halo (grey)}};
  \end{scope}
\end{tikzpicture}
   (body width, 120 mm)
\begin{tikzpicture}[x=1mm, y=1mm, font=\scriptsize\sffamily,
                    >={Stealth[length=1.4mm]}]
  \def\s{3.4}  % cell size
  % ---- left: full grid, four quadrants (rows counted from the top)
  \foreach \r in {1,...,8} {
    \foreach \c in {1,...,10} {
      \ifnum\r<5
        \ifnum\c<6 \def\col{orange!45}\else \def\col{yellow!60}\fi
      \else
        \ifnum\c<6 \def\col{blue!30}\else \def\col{green!45}\fi
      \fi
      \fill[\col, draw=white, line width=0.3pt] ({(\c-1)*\s},{-\r*\s}) rectangle ({\c*\s},{-(\r-1)*\s});
    }
  }
  \draw[ln-muted] (0,0) rectangle (10*\s,-8*\s);
  \draw[ln-muted, thick] (5*\s,0) -- (5*\s,-8*\s);
  \draw[ln-muted, thick] (0,-4*\s) -- (10*\s,-4*\s);
  \foreach \c/\r/\t in {5/5/C, 5/4/1, 4/5/2, 6/5/3, 5/6/4} {
    \node[font=\scriptsize\sffamily\bfseries] at ({(\c-0.5)*\s},{-(\r-0.5)*\s}) {\t};
  }
  \node[text=ln-muted] at (2.5*\s,1.8) {P0};
  \node[text=ln-muted] at (7.5*\s,1.8) {P1};
  \node[text=ln-muted] at (-2.5,-6*\s) {P2};
  \node[text=ln-muted] at (10*\s+2.5,-6*\s) {P3};
  \node[anchor=north, text=ln-muted, align=center] at (5*\s,-8*\s-1)
    {\iflangja{4プロセスに分割した格子}{grid split among four processes}};
  % ---- right: the blue region plus its halo
  \begin{scope}[shift={(52,0)}]
    \foreach \c in {1,...,5} {
      \fill[gray!45, draw=white, line width=0.3pt] ({(\c-1)*\s},{-4*\s}) rectangle ({\c*\s},{-3*\s});
    }
    \foreach \r in {5,...,8} {
      \fill[gray!45, draw=white, line width=0.3pt] ({5*\s},{-\r*\s}) rectangle ({6*\s},{-(\r-1)*\s});
      \foreach \c in {1,...,5} {
        \fill[blue!30, draw=white, line width=0.3pt] ({(\c-1)*\s},{-\r*\s}) rectangle ({\c*\s},{-(\r-1)*\s});
      }
    }
    \draw[ln-muted] (0,-4*\s) rectangle (5*\s,-8*\s);
    \foreach \c/\r/\t in {5/5/C, 5/4/1, 4/5/2, 6/5/3, 5/6/4} {
      \node[font=\scriptsize\sffamily\bfseries] at ({(\c-0.5)*\s},{-(\r-0.5)*\s}) {\t};
    }
    \draw[->, thick, ln-accent] (2.5*\s,-1.2*\s) -- (2.5*\s,-2.9*\s)
      node[pos=0, above, text=ln-accent] {\iflangja{P0から}{from P0}};
    \draw[->, thick, ln-accent] (8*\s,-6.5*\s) -- (6.2*\s,-6.5*\s)
      node[pos=0, right, text=ln-accent] {\iflangja{P3から}{from P3}};
    \node[anchor=north, text=ln-muted, align=center] at (3*\s,-8*\s-1)
      {\iflangja{P2の領域とハロー（灰色）}{P2's region and its halo (grey)}};
  \end{scope}
\end{tikzpicture}

  \caption{A stencil on a grid distributed over four processes.  Left:
    the element $C$ and its four neighbours straddle three regions.
    Right: process~2's region with the halo cells (grey) received from
    processes 0 and~3 before the update.}
  \label{fig:l02-stencil-halo}
\end{figure}

\begin{definition}[Halo and boundary exchange]
  The cells of a neighbouring region that a process reads but does not
  own form its \term{halo (ghost cells)}[ハロー領域]\index{ghost cells|see{halo (ghost cells)}}: for a stencil of
  radius~1, one row or column along each shared edge
  (\cref{fig:l02-stencil-halo}, right).  Before each update step the
  processes perform a \term{boundary exchange}[境界交換]: every process
  sends the edge cells its neighbours need and receives its halo from
  them, after which every stencil read is local.\margin{Each process sends
  and receives one message per neighbouring edge per step: with a
  $2\times2$ decomposition, two sends and two receives.  Doubling the
  stencil radius doubles the halo width, not the message count.}
\end{definition}

Whenever a data dependence crosses the boundary between processes, this
exchange is the price of distributed memory.  The same structure returns
in Lesson~3: a CUDA thread block owns a tile of the grid in its shared
memory and must load the halo around the tile before computing, except
that there the halo is read from global memory rather than received in a
message.\caution{A process may not read a halo cell before the message
carrying it has arrived, nor send its edge cells before it has updated
them: a distributed stencil alternates communication and computation,
and only the interior cells can be updated while the halo is still in
flight.}

\begin{keypoint}[Lesson summary]
  Parallel programs reuse a few patterns: master/worker with a task
  queue, loop parallelism, SPMD, fork/join and pipeline.  In shared
  memory, threads communicate through one address space and OpenMP
  annotates a serial program with directives for parallel loops, thread
  count, scheduling (static, dynamic, guided), data sharing (private,
  shared) and synchronisation (barrier, critical, atomic); an unprotected
  update of a shared variable is a data race.  In distributed memory,
  processes identified by rank exchange messages with MPI, point-to-point
  (\texttt{MPI\_Send}/\texttt{MPI\_Recv}) or collectively
  (\texttt{MPI\_Bcast}).  A distributed stencil needs a boundary exchange
  of halo cells before every step, the model for the halo loading of a
  CUDA tile in Lesson~3.
\end{keypoint}

\end{document}
